\documentclass{article}
\usepackage[margin=3cm]{geometry}

\usepackage{graphicx}

\usepackage[table]{xcolor}
\usepackage{tikz}
\usetikzlibrary{shapes,arrows}
\definecolor{xlightgray}{HTML}{D1DBE4}
\definecolor{xgreen}{HTML}{B9C89F}
\definecolor{xred}{HTML}{F19F9E}
\definecolor{xyellow}{HTML}{F6DA71}
\definecolor{xblue}{HTML}{ACE7EF}
\definecolor{xpurple}{HTML}{C3A4E7}
\usepackage{todonotes}
\usepackage{booktabs,tabularx,makecell,hyperref}
\definecolor{rowgray}{RGB}{246,246,246}

\usepackage{pdflscape}
\usepackage{environ}

\usetikzlibrary{arrows.meta,positioning,fit,calc,decorations.pathmorphing}

% Math
\usepackage{amsmath,amssymb}
\usepackage{mathtools}

% Algorithmic
\usepackage[noend]{algpseudocode}
\usepackage{amsthm}
\usepackage{cleveref}

\usepackage{float}
\floatstyle{ruled}
\newfloat{algo}{htbp}{algo}
\floatname{algo}{Algorithm}

% --- Custom commands ---
\newcommand{\genesis}{\ensuremath{B_{\mathrm{genesis}}}}
\newcommand{\T}[0]{\ensuremath{\mathcal{T}}}
\newcommand{\slot}[0]{\ensuremath{\operatorname{\mathrm{slot}}}}

\newcommand{\target}{\mathsf{target}}
\newcommand{\kind}{\mathsf{kind}}
\newcommand{\height}{\mathsf{height}}
\newcommand{\validator}{\mathsf{validator}}
\newcommand{\just}{\mathsf{justify}}
\newcommand{\notar}{\mathsf{notarize}}
\newcommand{\parent}{\mathsf{parent}}

% Algorithmic event environment
\algdef{SE}[EVENT]{Event}{EndEvent}[1]{\textbf{at}\ #1\ \algorithmicdo}{\algorithmicend\ \textbf{event}}
\algtext*{EndEvent}

\pagestyle{plain}

\renewcommand{\paragraph}[1]{\vspace{7pt}\noindent\textbf{#1}}

\algrenewcommand\textproc{}

\usepackage{boxedminipage}

% Theorem environments
\theoremstyle{plain}
\newtheorem{theorem}{Theorem}
\newtheorem*{theorem*}{Theorem}
\newtheorem{corollary}{Corollary}
\newtheorem{lemma}{Lemma}
\newtheorem{proposition}{Proposition}

\theoremstyle{definition}
\newtheorem{definition}[theorem]{Definition}

\theoremstyle{remark}
\newtheorem{remark}[theorem]{Remark}

\crefname{condition}{condition}{conditions}
\newtheorem{condition}{Condition}
\newtheorem{assumption}[theorem]{Assumption}
\newcommand{\GST}{\ensuremath{\mathsf{GST}}}


\title{Fresh Simplex with Prefix Notarizations}

\begin{document}
\maketitle

\section{Model and definitions}

We consider a system of~$n$ validators of which at most~$f$ are adversarial, where $n \geq 3f + 1$.
A \emph{quorum} is any set of at least $2n/3$ validators.
Any two quorums overlap in at least $n/3 + 1 > f$ validators, so their intersection contains at least one honest validator.

\begin{definition}[Height]\label{def:height}
\emph{Height} is a counter for the finality gadget, decoupled from slots, rounds, and epochs.
The genesis height is $0$.
Height advances by exactly~$1$ via one of two mechanisms: fixed-target justification (Definition~\ref{def:justification}) or prefix notarization (Definition~\ref{def:notarization}).
There are no timeouts.
\end{definition}

\begin{definition}[Block]\label{def:block}
A \emph{block} has a \emph{parent} (another block), a \emph{slot} field, and contains a set of votes (Definition~\ref{def:vote}).
The \emph{genesis block}~$\genesis$ has no parent and slot~$0$.
Slots are strictly increasing along any chain: for every non-genesis block~$B$, $B.\slot > B.\parent.\slot$.
\end{definition}

\begin{definition}[Chain and divergence]\label{def:chain}
A \emph{chain} is a sequence of blocks forming a path from the genesis block, where each block's parent is its predecessor in the sequence.
Since a chain is uniquely determined by its last block, we identify chains with their last block and use the two interchangeably.
Two chains \emph{conflict} if neither is a prefix of the other.
The \emph{divergence point} of two conflicting chains~$X$ and~$Y$ is the last block they share.
\end{definition}

Blocks form a tree rooted at~$\genesis$; each block has a unique path to the root.
We write $B \preceq B'$ (equivalently $B' \succeq B$) when $B$ is an ancestor of~$B'$ (or $B = B'$), and $B \sim B'$ when $B$ and $B'$ are \emph{compatible}: $B \preceq B'$ or $B' \preceq B$.
Among blocks on the same chain, \emph{higher} means further from genesis (i.e., $B' \succ B$).

\begin{definition}[Vote]\label{def:vote}
A \emph{vote} is a signed tuple
\[
  (\validator,\; \height,\; \slot,\; \target,\; \kind,\; \mathsf{finalize\_height},\; \mathsf{finalize\_target})
\]
where:
\begin{itemize}
  \item $\validator$ identifies the signer;
  \item $\height$ is a height;
  \item $\slot$ is a slot number used for tiebreaking in the fork-choice store;
  \item $\target$ is a block (never~$\bot$);
  \item $\kind \in \{\just,\, \notar\}$ distinguishes the vote type;
  \item $\mathsf{finalize\_height}$ is either a height or~$\bot$;
  \item $\mathsf{finalize\_target}$ is either a block or~$\bot$, and is $\bot$ exactly when $\mathsf{finalize\_height} = \bot$.
\end{itemize}
A vote with $\mathsf{finalize\_height} = h_f$ and $\mathsf{finalize\_target} = T_f$ explicitly commits the signer to finalizing~$T_f$ at height~$h_f$.
A vote included on a chain is verified to reference a target (and, when present, a finalize target) that actually exists on that chain.
\end{definition}

\paragraph{State.}
The \emph{state after block~$B$}, written $\sigma[B]$, is the tuple
\[
  (L,\; s,\; h,\; s_h,\; \mathsf{jtargets},\; \mathsf{ntargets},\; J,\; h_j,\; N,\; h_n,\; F,\; P).
\]
A block~$B'$ is \emph{on the chain ending at~$B$} iff $B' \preceq B$.
\begin{itemize}
  \item $L$: the \emph{chain head}, equal to~$B$ after processing block~$B$.
  \item $s$: the \emph{current slot}, incremented once per slot by $\textsf{processSlot}$ (Figure~\ref{alg:state-machine}).
  \item $h$: the \emph{current height}.
  \item $s_h$: the \emph{slot at which the current height began on this chain} --- the slot of the block whose processing last advanced~$h$ on this chain (genesis's slot~$0$ if no advance has yet occurred).
  \item $\mathsf{jtargets}[1..n]$: for each validator~$i$, the block~$i$ voted for with a $\just$ vote at the current height ($\bot$ if no such vote has been processed).
  \item $\mathsf{ntargets}[1..n]$: for each validator~$i$, the block~$i$ voted for with a $\notar$ vote at the current height ($\bot$ if no such vote has been processed).
  \item $J$: the \emph{justified block}, from the most recent justification event.
  \item $h_j$: the height at which~$J$ was justified.
  \item $N$: the \emph{latest notarized block}, from the most recent justification or notarization event ($J$ is automatically also a notarization, as noted in Definition~\ref{def:notarization}).
  \item $h_n$: the height at which~$N$ was notarized (or justified).
  \item $F$: the most recently finalized block.
  \item $P \subseteq \{1, \ldots, n\}$: the validators who have cast a vote with $\mathsf{finalize\_height} = h_j$ and $\mathsf{finalize\_target} = J$.
\end{itemize}
The genesis state is
\[
  \sigma[\genesis] = (\genesis,\; 0,\; 0,\; 0,\; \bot^n,\; \bot^n,\; \genesis,\; 0,\; \genesis,\; 0,\; \genesis,\; \emptyset).
\]
When unambiguous we drop the chain prefix and refer to components of $\sigma[L]$ simply as $L$, $s$, $h$, $s_h$, $\mathsf{jtargets}$, etc.

\begin{definition}[State-height of a block]\label{def:sheight}
The \emph{state-height} of a block~$B$ (on its own chain) is $\sigma[B].h$, i.e.\ the height of the finality gadget after processing~$B$ as the head of its own chain.
\end{definition}

\paragraph{Height freshness.}
Votes in this system are processed under a \emph{height-freshness} rule.
A vote~$v$ with target~$T$ at height~$v.\height$ is processed on a chain with current state $\sigma[L]$ only if
\[
  \sigma[L].h = v.\height \;\wedge\; T \preceq L \;\wedge\; T.\slot \geq \sigma[L].s_h.
\]
The first condition says the chain is currently at the right height; the second says the target lies on this chain; the third says $T$'s slot is no earlier than the slot at which the chain's current height began.
Intuitively, a fresh vote for target~$T$ at height~$h$ means ``the voter saw~$T$ as a chain of height~$h$''.

\begin{remark}[Freshness in terms of~$T$'s own chain]\label{rem:fresh-equiv}
The slot-based check is equivalent to the characterization ``$T$'s own chain reached state-height~$h$'', i.e., $\sigma[T].h = h$.
Reason: state-height is monotonically non-decreasing along any chain (Lemma~\ref{lem:slotmono}), it equals~$h$ on~$L$, and it advanced from~$h-1$ to~$h$ at the block whose slot is~$\sigma[L].s_h$.
So the set of blocks on~$L$'s chain at state-height exactly~$h$ is precisely $\{B \preceq L : B.\slot \geq \sigma[L].s_h\}$.
In particular, $T.\slot \geq \sigma[L].s_h$ and $T \preceq L$ together imply $\sigma[T].h = h$.
\end{remark}

This rule is central: it ensures every justification or notarization produced on a chain at height~$h$ is witnessed by a set of targets whose own chains reached exactly height~$h$.

\begin{definition}[Fixed-target justification]\label{def:justification}
A block~$T$ is \emph{justified at height~$h$} on a chain iff at least $2n/3$ validators have been recorded in $\mathsf{jtargets}$ with target~$T$:
\[
  \bigl|\{i \,:\, \mathsf{jtargets}[i] = T\}\bigr| \;\geq\; 2n/3.
\]
Only $\just$ votes whose target is~$T$ itself contribute; descendants of~$T$ do not count.
\end{definition}

\begin{definition}[Prefix notarization]\label{def:notarization}
A block~$T$ is \emph{notarized at height~$h$} on a chain ending at~$B$ iff walking from~$B$ toward genesis and summing the recorded $\notar$ counts, the cumulative total first reaches $2n/3$ at~$T$:
\[
  T \;=\; \max_{\preceq}\Bigl\{ B' \preceq B \,:\, \textstyle\sum_{B'' \,:\, B' \preceq B'' \preceq B} |\{i : \mathsf{ntargets}[i] = B''\}| \;\geq\; 2n/3 \Bigr\}
\]
where the maximum is taken with respect to~$\preceq$ (i.e., the highest such~$T$), when the set on the right is non-empty; if the set is empty, no block is notarized at this height on this chain, and the notarization branch of $\textsf{processHeight}$ does not fire.
Equivalently, whenever notarization does fire, $|\{i : T \preceq \mathsf{ntargets}[i] \preceq B\}| \geq 2n/3$.
Unlike justification, notarization uses \emph{prefix counting}: $\notar$ votes for descendants of~$T$ contribute.

A justification of~$T$ is also a notarization of~$T$: a justification requires $2n/3$ $\just$ votes for~$T$ specifically, which trivially satisfies the prefix-count threshold at~$T$.
\end{definition}

\begin{definition}[Finality]\label{def:finality}
The justified block~$J$ is \emph{finalized} when $|P| \geq 2n/3$.
We say $J$ is \emph{finalized at height~$h_j$}, where $h_j$ is the height at which~$J$ was justified.
\end{definition}

\begin{definition}[Slashing condition]\label{def:slashing}
A validator is \emph{slashable} (or an \emph{equivocator}) if it signed two votes~$a$, $b$ (possibly $a = b$) satisfying the \emph{target / finalize-target conflict} rule:
$b$ carries a finalize commitment ($b.\mathsf{finalize\_target} \neq \bot$), $a$ has been cast at $b$'s commit height, and $a$'s target disagrees with $b$'s finalize target,
\begin{align*}
  &a.\height = b.\mathsf{finalize\_height} \;\wedge\;
  a.\target \neq b.\mathsf{finalize\_target}.
\end{align*}
In words, a validator who signs $\mathsf{finalize\_target} = T_f$ at $\mathsf{finalize\_height} = h_f$ must not also sign \emph{any} vote (justification or notarization) with $\height = h_f$ and a different target.
We refer to this rule as~E1 in what follows.
\end{definition}

There is no direct slashing rule linking two $\just$ (or two $\notar$) votes at the same height; honest per-height uniqueness of $\just$ votes is instead imposed as a behavioral constraint on honest validators (Definition~\ref{def:voting-rule}).
A validator who never sets $\mathsf{finalize\_target} \neq \bot$ trivially avoids slashing.

\paragraph{State transitions.}
The state machine exposes two per-chain update routines (Figure~\ref{alg:state-machine}): $\textsf{processSlot}$ and $\textsf{processBlock}$.
A node invokes $\textsf{processSlot}$ once per slot on every chain it maintains, and $\textsf{processBlock}$ additionally whenever a block~$B$ is processed on that chain.

$\textsf{processSlot}(\sigma)$ advances the slot counter via $s \gets s + 1$ and then invokes $\textsf{processHeight}$ once.
$\textsf{processHeight}$ checks, in order: finality ($|P| \geq 2n/3 \Rightarrow F \gets J$); justification (if some $T$ has $|\{i : \mathsf{jtargets}[i] = T\}| \geq 2n/3$, set $(J, N) \gets (T, T)$, $(h_j, h_n) \gets (h, h)$, $P \gets \emptyset$, $h \gets h+1$, $s_h \gets L.\slot$, and reset the trackers); and finally prefix notarization (walk from~$L$ toward genesis summing $\mathsf{ntargets}$-counts and let $T^*$ be the highest block at which the cumulative total first reaches $2n/3$, then set $(N, h_n) \gets (T^*, h)$, $h \gets h+1$, $s_h \gets L.\slot$, and reset the trackers).
A single $\textsf{processHeight}$ invocation fires at most one height-advance: the justification and notarization branches each reset $\mathsf{jtargets}$ and $\mathsf{ntargets}$, so no further quorum can form without fresh votes being processed.

$\textsf{processBlock}(\sigma, B)$ is the block-driven update, with precondition $\sigma.s = B.\slot$: the caller advances $\sigma.s$ up to~$B.\slot$ via $\textsf{processSlot}$ first, then $\textsf{processBlock}$ sets $L \gets B$ and processes each vote~$v$ in~$B$ via $\textsf{processVote}$, which updates $\mathsf{jtargets}$, $\mathsf{ntargets}$, and~$P$.
On a fresh vote, $\mathsf{ntargets}[i]$ is updated to $\max_{\preceq}(\mathsf{ntargets}[i],\, v.\target)$ regardless of kind, so $\just$ votes also contribute to the prefix-notarization tally; additionally, if $v.\kind = \just$ then $\mathsf{jtargets}[i] \gets v.\target$.
The $\max_{\preceq}$ is well-defined because both arguments lie on $\sigma.L$ at state-height~$\sigma.h$ (treating $\bot$ as~$\genesis$).
The $P$-gate adds~$i$ to~$P$ iff $v.\mathsf{finalize\_height} = \sigma.h_j$ and $v.\mathsf{finalize\_target} = \sigma.J$; this uses only local $h_j$,~$J$, not freshness.
$\textsf{processBlock}$ does not itself advance~$h$: height advances are the responsibility of $\textsf{processSlot}$, which runs at every slot (including slots without a block).

\begin{figure}[H]
\caption{State machine}\label{alg:state-machine}
\begin{center}
\begin{boxedminipage}{\textwidth}
\begin{algorithmic}[1]

\Function{$\textsf{processSlot}$}{$\sigma$} \Comment{Runs once per slot on every chain, even without a block}
  \State $\sigma.s \gets \sigma.s + 1$
  \State $\sigma \gets \textsf{processHeight}(\sigma)$
  \State \Return $\sigma$
\EndFunction
\vspace{3pt}

\Function{$\textsf{processBlock}$}{$\sigma,\; B$}
  \State \textbf{assert} $\sigma.s = B.\slot$ \Comment{Precondition: caller advanced $\sigma.s$ via $\textsf{processSlot}$}
  \State $\sigma.L \gets B$
  \For{each vote $v$ in $B$} $\sigma \gets \textsf{processVote}(\sigma,\, v)$ \EndFor
  \State \Return $\sigma$
\EndFunction
\vspace{3pt}

\Function{$\textsf{processVote}$}{$\sigma,\; v$}
  \State $i \gets v.\validator$
  \State $\mathsf{fresh} \gets (v.\height = \sigma.h) \land (v.\target \preceq \sigma.L) \land (v.\target.\slot \geq \sigma.s_h)$
  \If{$\mathsf{fresh}$}
    \State $\sigma.\mathsf{ntargets}[i] \gets \max_{\preceq}(\sigma.\mathsf{ntargets}[i],\, v.\target)$ \Comment{both kinds bump $\mathsf{ntargets}$}
    \If{$v.\kind = \just$}
      \State $\sigma.\mathsf{jtargets}[i] \gets v.\target$
    \EndIf
  \EndIf
  \If{$v.\mathsf{finalize\_height} = \sigma.h_j$ \textbf{and} $v.\mathsf{finalize\_target} = \sigma.J$}
    \State $\sigma.P \gets \sigma.P \cup \{i\}$
  \EndIf
  \State \Return $\sigma$
\EndFunction
\vspace{3pt}

\Function{$\textsf{processHeight}$}{$\sigma$}
  \If{$|\sigma.P| \geq 2n/3$} \Comment{Finality}
    \State $\sigma.F \gets \sigma.J$
  \EndIf
  \State $T \gets \arg\max_{T' \in \{\sigma.\mathsf{jtargets}[i] \,:\, \sigma.\mathsf{jtargets}[i] \neq \bot\}} |\{i : \sigma.\mathsf{jtargets}[i] = T'\}|$ \Comment{$T \gets \bot$ if the set is empty; the next guard then skips this branch via the $2n/3$ check}
  \If{$T \neq \bot$ \textbf{and} $|\{i : \sigma.\mathsf{jtargets}[i] = T\}| \geq 2n/3$} \Comment{Justification (also a notarization)}
    \State $\sigma.J \gets T$, \; $\sigma.h_j \gets \sigma.h$, \; $\sigma.P \gets \emptyset$
    \State $\sigma.N \gets T$, \; $\sigma.h_n \gets \sigma.h$
    \State $\sigma.h \gets \sigma.h + 1$, \; $\sigma.s_h \gets \sigma.L.\slot$
    \State reset $\mathsf{jtargets}, \mathsf{ntargets}$
    \State \Return $\sigma$
  \EndIf
  \State $\mathsf{count} \gets 0$ \Comment{Prefix notarization walk}
  \For{$B$ from $\sigma.L$ down to $\genesis$}
    \State $\mathsf{count} \gets \mathsf{count} + |\{i : \sigma.\mathsf{ntargets}[i] = B\}|$
    \If{$\mathsf{count} \geq 2n/3$} \Comment{Notarization fires at $B$}
      \State $\sigma.N \gets B$, \; $\sigma.h_n \gets \sigma.h$
      \State $\sigma.h \gets \sigma.h + 1$, \; $\sigma.s_h \gets \sigma.L.\slot$
      \State reset $\mathsf{jtargets}, \mathsf{ntargets}$
      \State \Return $\sigma$
    \EndIf
  \EndFor
  \State \Return $\sigma$
\EndFunction
\vspace{3pt}

\Function{$\textsf{finalizeConflict}$}{$v_1, v_2$} \Comment{Does $v_1$ conflict with $v_2$'s finalize commitment?}
  \State \Return $v_2.\mathsf{finalize\_target} \neq \bot \;\wedge\; v_1.\height = v_2.\mathsf{finalize\_height} \;\wedge\; v_1.\target \neq v_2.\mathsf{finalize\_target}$
\EndFunction
\vspace{3pt}

\Function{$\textsf{isSlashable}$}{$v_1, v_2$}
  \State \Return $v_1.\validator = v_2.\validator$
  \Statex \hspace{3em}$\land \;(\textsf{finalizeConflict}(v_1, v_2) \lor \textsf{finalizeConflict}(v_2, v_1))$
\EndFunction

\end{algorithmic}
\end{boxedminipage}
\end{center}
\end{figure}

\paragraph{Freshness as a structural invariant.}
If $v$ is recorded on chain~$L$ at state-height~$h$, then $v.\target \preceq L$ with $\sigma[v.\target].h = h$ (Remark~\ref{rem:fresh-equiv}).
Consequently any block justified or notarized at height~$h$ on some chain has $\geq 2n/3$ witnessing targets, each at state-height~$h$ on its own chain.

\section{Accountable safety}

\begin{lemma}[Height progression]\label{lem:heightprog}
$\textsf{processHeight}$ increments $h$ by at most~$1$ per invocation, and each increment is gated by a $\geq 2n/3$ justification or prefix-notarization quorum.
In particular, any chain reaching height~$h'$ from~$h < h'$ has advanced through each intermediate height via such a quorum.
\end{lemma}

\begin{proof}
The only assignments to~$h$ are $h \gets h + 1$ in the justification and notarization branches of $\textsf{processHeight}$, each guarded by its $\geq 2n/3$ quorum.
Reaching $h'$ from $h$ therefore requires $h' - h$ successive increments, each so gated.
\end{proof}

\begin{lemma}[Checkpoint monotonicity]\label{lem:slotmono}
On any chain, the fields $s$, $h$, $h_j$, $h_n$, $s_h$, $F$, $J$, $N$ are non-decreasing along~$\preceq$, with $F \preceq J \preceq N$ at all times.
Moreover $J$ and $h_j$ advance only on justification events, $h_n$ advances only on height-advances, and $s_h$ advances only on height-advances, each time being set to the slot of the advancing block.
Consequently, whenever $\sigma[Y].h = h$, the blocks on~$Y$ at state-height exactly~$h$ are $\{B' \preceq Y : B'.\slot \geq \sigma[Y].s_h\}$.
\end{lemma}

\begin{proof}
We check each field in turn, reading off the assignments in Figure~\ref{alg:state-machine}.

\emph{Scalar fields.}
$s$ only increases (via $s \gets s + 1$ in $\textsf{processSlot}$).
$h$ only increases (by~$1$, in the justification or notarization branch of $\textsf{processHeight}$).
$h_j$ and $h_n$ are set to the current~$h$ at each $J$- or $N$-update, so they too are non-decreasing.
$s_h$ is updated only in the height-advance branches, jointly with $h \gets h+1$, and set to $L.\slot$; since block slots strictly increase along~$\preceq$ (Definition~\ref{def:block}), $s_h$ is non-decreasing.

\emph{Block fields $F$, $J$, $N$.}
$F$ is updated only via $F \gets J$, so $F \preceq J$ follows once $J$-monotonicity is established.
Consider a height-advance on a chain~$Y$ with pre-state height~$h$, firing at some target~$T$:
\begin{itemize}
  \item By Remark~\ref{rem:fresh-equiv}, $\sigma[T].h = h$.
  For a justification, each contributing $\just$ vote has target~$T$ with slot $\geq s_h$ (by freshness).
  For a prefix notarization, every recorded entry of $\mathsf{ntargets}$ on~$Y$ points at a target with slot $\geq s_h$ (again by freshness), so all contributors lie in the height-$h$ slot interval; the cumulative prefix count can only first reach $2n/3$ inside that interval, so the firing block~$T$ satisfies $T.\slot \geq s_h$.
  \item The previous value $N_{\mathrm{old}}$ has $\sigma[N_{\mathrm{old}}].h \leq h-1$.
  Combined with $\sigma[T].h = h$ and $T, N_{\mathrm{old}} \preceq Y$, state-height monotonicity along~$\preceq$ on~$Y$ yields $T \succ N_{\mathrm{old}}$.
\end{itemize}
The justification branch sets $J = N = T$ jointly and the notarization branch sets only $N \gets T$, leaving $J$ unchanged.
Thus $N$ strictly advances at every height-advance; $J$ strictly advances at every justification; and $J \preceq N$ is preserved, as is $J$-monotonicity.

\emph{Height-$h$ interval.}
When $\sigma[Y].h = h$, the current value of $s_h$ was last set to the slot of the block that advanced state-height to~$h$ along~$Y$.
By strict slot increase along~$\preceq$ (Definition~\ref{def:block}) and state-height monotonicity, the ancestors $B' \preceq Y$ with $\sigma[B'].h = h$ are exactly those with $B'.\slot \geq s_h$.
\end{proof}

\begin{lemma}[Any chain past a finalized height contains the finalized block]\label{lem:mainsafety}
Unless $\geq n/3$ validators are slashable: if~$C$ is finalized at height~$h$, then $C \preceq B$ for every block~$B$ with $\sigma[B].h > h$.
\end{lemma}

\begin{proof}
Since $\sigma[B].h > h$, Lemma~\ref{lem:heightprog} gives a block $B^* \preceq B$ at which the advance from height~$h$ to~$h+1$ fired, via a height-$h$ justification or prefix-notarization quorum~$Q$.
Each $i \in Q$ contributed a vote~$a_i$ with $a_i.\height = h$ processed freshly on the chain ending at~$B^*$ (directly for $\just$; for prefix $\notar$, the $\mathsf{ntargets}$-update rule records only fresh votes).
Finality of~$C$ at height~$h$ gives a quorum~$Q_F$ of validators each signing a vote~$b$ with $b.\mathsf{finalize\_height} = h$ and $b.\mathsf{finalize\_target} = C$.

By quorum intersection $|Q \cap Q_F| \geq n/3$; pick any $v \in Q \cap Q_F$ not slashable.
Applying E1 (Definition~\ref{def:slashing}) to $v$'s votes~$a \in Q$ and $b \in Q_F$: $b$ carries a finalize commitment to~$C$ at height~$h$, and $a.\height = h$, so $a.\target \neq C$ would be slashable.
Hence $a.\target = C$, and freshness gives $a.\target \preceq B^*$, so $C \preceq B^* \preceq B$.
\end{proof}

\begin{lemma}[Finalized blocks form a chain]\label{lem:finchain}
Unless $\geq n/3$ validators are slashable: for any two finalized checkpoints $(C, h)$ and $(C', h')$, blocks~$C$ and~$C'$ are compatible; in particular, if $h \leq h'$ then $C \preceq C'$.
\end{lemma}

\begin{proof}
Relabelling if necessary, assume $h \leq h'$.

\emph{Case $h = h'$.}
Finality of~$C$ at height~$h$ gives a quorum~$Q_F$ each signing a vote with $\mathsf{finalize\_height} = h$ and $\mathsf{finalize\_target} = C$.
Finality of~$C'$ at height~$h$ implies justification of~$C'$ at height~$h$, giving a $\just$-quorum~$Q'$ each signing a vote with $\height = h$ and $\target = C'$.
By quorum intersection $|Q_F \cap Q'| \geq n/3$; since $< n/3$ validators are slashable, pick $v \in Q_F \cap Q'$ not slashable.
E1 (Definition~\ref{def:slashing}) applied to $v$'s votes in~$Q_F$ and~$Q'$ forces $C = C'$.

\emph{Case $h < h'$.}
Let~$Z$ be the chain (identified with its tip) on which~$C'$ was finalized, taken at the post-state where the finalize-quorum for~$C'$ closed.
At that post-state, $\sigma[Z].h \geq h' > h$, with $C' \preceq Z$ and $\sigma[C'].h = h'$ (Remark~\ref{rem:fresh-equiv}).
Applying Lemma~\ref{lem:mainsafety} at~$Z$ yields $C \preceq Z$, and Remark~\ref{rem:fresh-equiv} applied on~$C$'s own finalizing chain gives $\sigma[C].h = h$.
Hence $C$ and~$C'$ both lie on~$Z$, so they are comparable under~$\preceq$; state-height monotonicity along~$\preceq$ on~$Z$ (Lemma~\ref{lem:slotmono}) combined with $\sigma[C].h = h < h' = \sigma[C'].h$ forces $C \prec C'$.
\end{proof}

\begin{theorem}[Accountable safety]\label{thm:safety}
Unless $\geq n/3$ validators are slashable, no two conflicting blocks can be finalized.
\end{theorem}

\begin{proof}
Immediate from Lemma~\ref{lem:finchain}: any two finalized blocks are compatible.
\end{proof}

\begin{remark}[Ancestor notarizations are permitted]\label{rem:notarization-safety}
Because $\just$ votes also contribute to $\mathsf{ntargets}$ (Figure~\ref{alg:state-machine}), the prefix-notarization endpoint at height~$h$ on a chain containing a finalized block~$C$ may sit at some ancestor $T^* \preceq C$ on that chain.
This is consistent with safety, but means the per-chain invariants never rely on a strict $N \succeq C$ relation; the fork-choice store (Section~\ref{sec:store}) uses a greedy walk to handle this.
\end{remark}

\section{Fork-choice store}\label{sec:store}

A node may track multiple chains simultaneously.
The \emph{store} is the node-local data structure that mediates which chain the node follows.
It carries two sets of checkpoint records: a primary set $C_N$ of \emph{notarized checkpoint records}, one per height advance (justification or prefix notarization), and a sub-layer $C_J \subseteq C_N$ of \emph{justified checkpoint records} --- those records whose block was justified on some chain.
The store's root $R_s$ is computed by a greedy walk from the store-finalized block~$F_s$, following the highest-height notarized record that strictly descends from the current walk position, with $C_J$-membership as a tiebreak at equal heights.
The walk guarantees $R_s \succeq F_s$ even when records for strict ancestors of finalized blocks remain in~$C_N$ (Remark~\ref{rem:notarization-safety}).

\begin{definition}[Store]\label{def:store}
A node maintains a \emph{store}
\[
  \Sigma \;=\; (\sigma,\; R_s,\; F_s,\; \T,\; C_J,\; C_N),
\]
where:
\begin{itemize}
  \item $\sigma$ is the \emph{state map}, assigning each accepted block its per-chain post-state (Section~1);
  \item $R_s$ is the \emph{store root}, used as the root of the fork-choice subtree;
  \item $F_s$ is the \emph{store-finalized block};
  \item $\T$ is the \emph{block tree}, the set of accepted blocks;
  \item $C_N$ is the set of \emph{notarized checkpoint records} $(X, h, s)$, one per post-state of an accepted block: $X = \sigma[B].N$, $h = \sigma[B].h_n$, $s = X.\slot$;
  \item $C_J \subseteq C_N$ is the set of \emph{justified checkpoint records}: those records whose block was justified (not merely prefix-notarized) at height~$h$ on some chain.
\end{itemize}
The capital~$\Sigma$ distinguishes the store from the per-chain state~$\sigma$: $\sigma[B]$ is a single post-state, while $\Sigma$ aggregates the state map together with the fork-choice bookkeeping $(R_s, F_s, \T, C_J, C_N)$.
Where unambiguous we write $\Sigma.R_s$, $\Sigma.F_s$, $\Sigma.\T$, $\Sigma.C_J$, $\Sigma.C_N$, $\Sigma.\sigma$ for the corresponding fields.
The genesis store has $R_s = F_s = \genesis$, $\T = \{\genesis\}$, $C_J = C_N = \{(\genesis, 0, 0)\}$.
\end{definition}

\begin{remark}[Store-finalized invariant]\label{rem:fs-invariant}
At all times, $F_s$ is either~$\genesis$ or a block that was finalized on some chain, and $(F_s, h_{F_s}, F_s.\slot) \in C_J$ where $h_{F_s}$ is the height at which~$F_s$ was finalized.
Indeed, $F_s$ is only updated by $\textsf{updateFinalized}$ (Figure~\ref{alg:store}) to a value~$F_B = \sigma[B].F$, and $\sigma[B].F$ is set only by the finality branch of $\textsf{processHeight}$.
The $C_J$-membership holds by initialization for $F_s = \genesis$; for an update $F_s \gets F_B$, some ancestor~$B'$ of~$B$ (possibly~$B$ itself) had post-state $(J_{B'}, h_j^{B'}) = (F_B, h_{F_s})$, and its processing inserted $(F_B, h_{F_s}, F_B.\slot)$ into~$C_J$.
\end{remark}

The store's block-acceptance routine $\textsf{onBlock}$ (Figure~\ref{alg:store}) admits a block~$B$ only if $B.\parent \in \Sigma.\T$ and $\Sigma.F_s \preceq B$.
It runs $\textsf{processSlot}$ once for each slot from $\Sigma.\sigma[B.\parent].s + 1$ up to $B.\slot$ (so any height advance enabled by the accumulated tally fires before~$B$'s votes are folded in), then $\textsf{processBlock}(\cdot, B)$; the resulting post-state becomes $\Sigma.\sigma[B]$.
The tuples $(N_B, h_n^B, N_B.\slot)$ and $(J_B, h_j^B, J_B.\slot)$ (subscripts denoting fields of $\Sigma.\sigma[B]$) are inserted into $C_N$ and $C_J$ respectively; a justification fires $(J, N) \gets (T, T)$ and $(h_j, h_n) \gets (h, h)$ jointly, so the inserted tuples coincide and $C_J \subseteq C_N$ is preserved.
$\Sigma.R_s$ is then recomputed by the greedy walk $\textsf{computeRoot}$: starting from~$\Sigma.F_s$, each step picks the tuple $(X, h, s) \in \Sigma.C_N$ with $X$ a strict descendant of the current walk position maximizing the lex key $\bigl(h,\; \mathbf{1}[(X, h, s) \in \Sigma.C_J],\; s,\; \mathsf{hash}(X)\bigr)$ and moves to~$X$; the argmax is unique (block hashes are unique) and the walk terminates in at most $|\Sigma.C_N|$ steps.
The indicator bit gives a $C_J$-tuple a strict edge over any $C_N \setminus C_J$-tuple at the same height.
Finally, $\textsf{updateFinalized}$ sets $\Sigma.F_s \gets F_B$ whenever $F_B \succ \Sigma.F_s$ and $F_B \preceq \Sigma.R_s$.

\begin{figure}[H]
\caption{Store and fork-choice root}\label{alg:store}
\begin{center}
\begin{boxedminipage}{\textwidth}
\begin{algorithmic}[1]

\State \textbf{Genesis store} $\Sigma$: \;$\Sigma.R_s \gets \genesis$, \; $\Sigma.F_s \gets \genesis$, \; $\Sigma.\T \gets \{\genesis\}$,
\State \hspace{8.1em}$\Sigma.C_J \gets \{(\genesis, 0, 0)\}$, \; $\Sigma.C_N \gets \{(\genesis, 0, 0)\}$
\State \hspace{8.1em}$\Sigma.\sigma[\genesis] \gets (\genesis,\, 0,\, 0,\, 0,\, \bot^n,\, \bot^n,\, \genesis,\, 0,\, \genesis,\, 0,\, \genesis,\, \emptyset)$
\vspace{3pt}

\Function{$\textsf{onBlock}$}{$\Sigma,\; B$}
  \State \textbf{assert} $B.\parent \in \Sigma.\T$
  \State \textbf{assert} $\Sigma.F_s \preceq B$
  \State $\Sigma.\T \gets \Sigma.\T \cup \{B\}$
  \State $\sigma' \gets \Sigma.\sigma[B.\parent]$
  \While{$\sigma'.s < B.\slot$} \Comment{Advance through every intervening slot}
    \State $\sigma' \gets \textsf{processSlot}(\sigma')$
  \EndWhile
  \State $\Sigma.\sigma[B] \gets \textsf{processBlock}(\sigma',\; B)$
  \State $\Sigma.C_J \gets \Sigma.C_J \cup \{(\Sigma.\sigma[B].J,\; \Sigma.\sigma[B].h_j,\; \Sigma.\sigma[B].J.\slot)\}$
  \State $\Sigma.C_N \gets \Sigma.C_N \cup \{(\Sigma.\sigma[B].N,\; \Sigma.\sigma[B].h_n,\; \Sigma.\sigma[B].N.\slot)\}$
  \State $\Sigma.R_s \gets \textsf{computeRoot}(\Sigma)$
  \State $\Sigma \gets \textsf{updateFinalized}(\Sigma,\; \Sigma.\sigma[B].F)$
  \State \Return $\Sigma$
\EndFunction
\vspace{3pt}

\Function{$\textsf{computeRoot}$}{$\Sigma$} \Comment{Greedy walk from $\Sigma.F_s$}
  \State $X^* \gets \Sigma.F_s$
  \While{$\exists\, (X, h, s) \in \Sigma.C_N$ with $X \succ X^*$}
    \State $(X^*, \_, \_) \gets \arg\max_{(X,h,s) \in \Sigma.C_N,\; X \succ X^*} \bigl(h,\; \mathbf{1}[(X,h,s) \in \Sigma.C_J],\; s,\; \mathsf{hash}(X)\bigr)$
  \EndWhile
  \State \Return $X^*$
\EndFunction
\vspace{3pt}

\Function{$\textsf{updateFinalized}$}{$\Sigma,\; F_B$}
  \If{$F_B \succ \Sigma.F_s$ and $F_B \preceq \Sigma.R_s$}
    \State $\Sigma.F_s \gets F_B$
  \EndIf
  \State \Return $\Sigma$
\EndFunction
\vspace{3pt}

\Function{$\textsf{getHead}$}{$\Sigma$} \Comment{Deterministic choice under a global $\textsf{getHead}$ policy (e.g., highest-slot leaf)}
  \State \Return any block in~$\Sigma.\T$ that descends from~$\Sigma.R_s$
\EndFunction

\end{algorithmic}
\end{boxedminipage}
\end{center}
\end{figure}

Whenever $R_s \neq F_s$, $R_s$ is the last-visited block of the walk, so $R_s$ is the block of some record $(R_s, R_s.h, R_s.\slot) \in C_N$.

\paragraph{Leaf-only variant.}
Define $C_N^{\mathrm{leaf}} \subseteq \Sigma.C_N$ to be the subset of tuples $(X, h, s)$ contributed by the post-state of a \emph{leaf} block of~$\Sigma.\T$ (a block with no processed descendant), and analogously $C_J^{\mathrm{leaf}} \subseteq \Sigma.C_J$.
Let $\textsf{computeRootLeaf}(\Sigma)$ be $\textsf{computeRoot}$ run over $(C_N^{\mathrm{leaf}}, C_J^{\mathrm{leaf}})$.

\begin{theorem}[Leaf-only equivalence]\label{thm:leafequiv}
$\textsf{computeRootLeaf}(\Sigma)$ returns the same block as $\textsf{computeRoot}(\Sigma)$.
\end{theorem}

\begin{proof}
Let $\tau_i = (X_i, h_i, X_i.\slot) \in C_N$ be the tuple contributed by a non-leaf processed block~$B_i$.
Since $\T$ is finite, some processed leaf~$B_\ell$ descends from~$B_i$; let $\tau_\ell = (X_\ell, h_\ell, X_\ell.\slot) \in C_N^{\mathrm{leaf}}$ be its tuple.
By Lemma~\ref{lem:slotmono} along the chain from~$B_i$ to~$B_\ell$, both $h_n$ and~$N$ are non-decreasing, so $h_\ell \geq h_i$ and $X_\ell \succeq X_i$.
Moreover, $N$ advances only jointly with $h_n \gets h_n + 1$ (Lemma~\ref{lem:slotmono}), so $h_\ell = h_i$ forces $X_\ell = X_i$, whence $\tau_\ell = \tau_i \in C_N^{\mathrm{leaf}}$.
Thus every non-leaf tuple is either itself a leaf tuple, or its leaf replacement~$\tau_\ell$ has strictly greater height.

Run both walks in lockstep from $X^* = F_s$.
At each step, let $\tau^\star$ be the full walk's lex-max candidate over~$C_N$.
If $\tau^\star$ is a leaf tuple, the leaf walk picks the same candidate.
Otherwise $\tau^\star = \tau_i$ for some non-leaf~$B_i$, and $\tau_\ell$ is also a strict-descendant candidate: by lex-maximality of $\tau^\star$, $\tau_\ell$'s key is at most $\tau_i$'s, which rules out $h_\ell > h_i$; hence $h_\ell = h_i$, $\tau_\ell = \tau_i \in C_N^{\mathrm{leaf}}$, and $\tau^\star$ is already a leaf tuple.
The two walks make the same choice at each step, so they terminate at the same block.
\end{proof}

The leaf-only variant bounds the walk's candidate pool by the number of leaves of~$\T$.

\subsection{Store invariants and safety}

\subsubsection*{Unconditional invariants}

\begin{theorem}[$F_s \preceq R_s$ unconditionally]\label{thm:fleqr}
The store maintains $F_s \preceq R_s$ at all times.
\end{theorem}

\begin{proof}
Initially $F_s = R_s = \genesis$.
$\textsf{computeRoot}$ starts at $X^* = F_s$ and only moves to strict descendants, so $R_s \succeq F_s$ holds immediately after; and $\textsf{updateFinalized}$ only sets $F_s$ to values $F_B \preceq R_s$ (explicit guard).
\end{proof}

\begin{theorem}[Local irreversibility of finality]\label{thm:finperm}
Once a node sets $F_s = F$, $F_s$ descends from~$F$ at all future times.
\end{theorem}

\begin{proof}
$F_s$ is updated only by $\textsf{updateFinalized}$, which requires $F_B \succ F_s$.
\end{proof}

\begin{theorem}[Fork-choice consistency]\label{thm:fcconsistency}
Once a node sets $F_s = F$, $\textsf{getHead}$ returns a block descending from~$F$ at all future times.
\end{theorem}

\begin{proof}
By Theorems~\ref{thm:finperm} and~\ref{thm:fleqr}, $F \preceq F_s \preceq R_s$, and $\textsf{getHead}$ returns a descendant of~$R_s$.
\end{proof}

\subsubsection*{Further properties under $f < n/3$}

The key tool for the remaining store theorems is Lemma~\ref{lem:certchain}: every notarized checkpoint record whose height is at least that of a finalized block~$C$ is compatible with~$C$, and strictly descends from~$C$ when its height strictly exceeds that of~$C$.

\begin{lemma}[Checkpoint chain]\label{lem:certchain}
Unless $\geq n/3$ validators are slashable: if~$C$ is finalized at height~$h_F$ and $(X, h', s) \in C_N$ with $h' \geq h_F$, then $C$ and~$X$ are compatible; moreover, if $h' > h_F$ then $C \prec X$.
\end{lemma}

\begin{proof}
If $h_F = 0$ then $C = \genesis$ and compatibility is trivial; assume $h_F \geq 1$, so $(X, h', s)$ is not the initial genesis tuple.
The tuple was therefore inserted from the post-state of some processed block~$B$ with $(\sigma[B].N, \sigma[B].h_n) = (X, h')$, so $X \preceq B$.
Every $h_n$-update is paired with $\sigma.h \gets \sigma.h + 1$ (Lemma~\ref{lem:slotmono}), hence $\sigma[B].h \geq h' + 1 > h_F$; Lemma~\ref{lem:mainsafety} gives $C \preceq B$, and with $X \preceq B$ this yields $C \sim X$.

\emph{Strict descent when $h' > h_F$.}
By Remark~\ref{rem:fresh-equiv}, $\sigma[X].h = h'$ and $\sigma[C].h = h_F$.
Both $C, X$ lie on the chain ending at~$B$, so they are comparable; state-height monotonicity along~$\preceq$ (Lemma~\ref{lem:slotmono}) with $\sigma[C].h = h_F < h' = \sigma[X].h$ forces $C \prec X$.
\end{proof}

\begin{lemma}[Upgrade property of the walk]\label{lem:upgrade}
Unless $\geq n/3$ validators are slashable: if~$F$ is finalized at height~$h_F$, $(F, h_F, F.\slot) \in C_J$, and $F_s \prec F$, then $\textsf{computeRoot}$ produces $R_s \succeq F$.
\end{lemma}

\begin{proof}
The walk starts at $X^* = F_s \prec F$.
Since $(F, h_F, F.\slot) \in C_J \subseteq C_N$ is a strict-descendant candidate, the first step picks a tuple $(X, h, s)$ with key $\geq (h_F, 1, F.\slot, \mathsf{hash}(F))$; in particular $h \geq h_F$.

If $h > h_F$, Lemma~\ref{lem:certchain} gives $F \prec X$.

If $h = h_F$, the key's bit coordinate must be~$1$, so $(X, h_F, s) \in C_J$ and the full key equals $(h_F, 1, X.\slot, \mathsf{hash}(X))$; hence $X.\slot \geq F.\slot$.
Lemma~\ref{lem:certchain} gives $F \sim X$, so $F \preceq X$ or $X \prec F$; the latter would force $X.\slot < F.\slot$ on $F$'s chain (Definition~\ref{def:block}), contradicting $X.\slot \geq F.\slot$.
Hence $F \preceq X$.

Either way $F \preceq X$, and subsequent steps move to strict descendants of~$X$, so $R_s \succeq F$.
\end{proof}

\begin{theorem}[Local acceptance of finality updates]\label{thm:finlive}
Unless $\geq n/3$ validators are slashable: if a block~$B$ is processed by $\textsf{onBlock}$ and its post-state has finalized block~$F_B$, then after processing $F_s \succeq F_B$.
\end{theorem}

\begin{proof}
If $F_s \succeq F_B$ already, the guard $F_B \succ F_s$ fails and the invariant is maintained.
Otherwise, $F_s$ is genesis or a finalized block (Remark~\ref{rem:fs-invariant}), so Lemma~\ref{lem:finchain} gives $F_s \sim F_B$, and with $F_s \not\succeq F_B$ this forces $F_s \prec F_B$.
It remains to show $F_B \preceq R_s$ after $\textsf{computeRoot}$, so that the $\textsf{updateFinalized}$ guard passes.

Let~$h_F$ be the height at which~$F_B$ was finalized.
Since $F_B$ is finalized in~$B$'s post-state, some ancestor~$B'$ of~$B$ (possibly~$B$ itself) had post-state $(J_{B'}, h_j^{B'}) = (F_B, h_F)$; blocks are processed parent-first, so $B'$ has already been processed and $(F_B, h_F, F_B.\slot) \in C_J$.
Lemma~\ref{lem:upgrade} then gives $R_s \succeq F_B$.
\end{proof}

\begin{theorem}[Lock-in]\label{thm:lockin}
Unless $\geq n/3$ validators are slashable: if block~$F$ is finalized at height~$h$ on any chain, and some block whose post-state has $(F, h)$ as its justified checkpoint has been processed by the node, then $R_s \succeq F$ at all future times, and $\textsf{getHead}$ always returns a descendant of~$F$.
\end{theorem}

\begin{proof}
Processing the block inserts $(F, h, F.\slot)$ into~$C_J$, and $C_J$ is monotone, so this tuple persists.
Fix any future moment.
If $F \preceq F_s$, Theorem~\ref{thm:fleqr} gives $R_s \succeq F_s \succeq F$.
Otherwise, Remark~\ref{rem:fs-invariant} and Lemma~\ref{lem:finchain} give $F_s \sim F$, whence $F \npreceq F_s$ forces $F_s \prec F$; Lemma~\ref{lem:upgrade} then gives $R_s \succeq F$.
Either way $\textsf{getHead}$ (descending from~$R_s$) returns a descendant of~$F$.
\end{proof}

\begin{theorem}[Order independence]\label{thm:orderindep}
Unless $\geq n/3$ validators are slashable: the store state after processing a set of blocks depends only on which blocks have been processed, not on their order.
\end{theorem}

\begin{proof}
Each block's post-state is a deterministic function of itself and its parent's post-state, so $\sigma$ depends only on the processed set.
$C_N$ and $C_J$ are inserted-only, and tuple insertion is idempotent; both are recomputed deterministically from~$\sigma$, and hence depend only on the processed set.

For $F_s$: by Lemma~\ref{lem:finchain}, the values $\{\sigma[B].F : B \in \T\}$ form a chain with maximum $F_{\max}$, attained at some processed block~$B^\ast$.
$F_s$ only takes values from this chain (Remark~\ref{rem:fs-invariant}), so $F_s \preceq F_{\max}$ at all times; and once $B^\ast$ has been processed, Theorem~\ref{thm:finlive} gives $F_s \succeq F_{\max}$.
Hence $F_s$ stabilizes at $F_{\max}$ regardless of order.

Finally, $\textsf{computeRoot}$ is a deterministic function of $(F_s, C_N, C_J)$: distinct hashes make the key total order, the argmax at each step is unique, and the walk terminates in at most $|C_N|$ steps.
Thus $R_s$ depends only on the processed set.
\end{proof}

\begin{corollary}[Cross-node agreement]\label{cor:agreement}
Unless $\geq n/3$ validators are slashable: two nodes that have processed the same set of blocks have the same $(R_s, F_s, C_J, C_N)$.
\end{corollary}

\begin{proof}
Immediate from Theorem~\ref{thm:orderindep}.
\end{proof}

\section{Inactivity leak}

Write $X = \textsf{getHead}(\Sigma)$ for the \emph{canonical chain}.
The leak section extends the per-chain state of Section~1 (Figure~\ref{alg:state-machine}) with a penalty counter and a new transition function $\textsf{processInactivityPenalties}$, invoked each slot by an extended $\textsf{processSlot}$ (Figure~\ref{alg:leak-processslot}).
The extension is active in this section only; $\textsf{processBlock}$, $\textsf{processHeight}$, and $\textsf{processVote}$ are unchanged.

Two results follow.
Tightness (Theorem~\ref{thm:tightness}) is purely state-logic: it reads $\sigma.\mathsf{penalties}$ on any single chain and bounds its cumulative increment over periods of non-finality on that chain, assuming nothing about message delivery, honesty, or the size of the slashable set.
Fairness (Theorem~\ref{thm:fairness}) is local to a single validator: the vote an honest~$i$ produces per the honest rule (Definition~\ref{def:voting-rule}), once included on any extension of~$i$'s canonical chain still at the same state-height, exempts~$i$ from Layer~1 penalties at that state-height.
Layer-2 exemption is out of scope: it depends on timely quorum formation on a common target, not on any local property of~$i$.
Height advances reset $\mathsf{jtargets}$ and $\mathsf{ntargets}$, so the next slot's Layer-2 state does not depend on~$i$'s votes at prior heights and no ``lock'' issue arises.

\subsection{Penalty units}

\paragraph{State extension.}
Augment the per-chain state tuple with one field
\[
  \sigma.\mathsf{penalties}[1..n],
\]
a per-validator non-negative counter, initialized to~$0$ for every~$i$ at genesis.
No other field is added or modified.

\paragraph{Extended $\textsf{processSlot}$.}
Figure~\ref{alg:leak-processslot} replaces $\textsf{processSlot}$ from Figure~\ref{alg:state-machine} with a version that additionally calls $\textsf{processInactivityPenalties}$ on the pre- and post-height-advance snapshots.
The store of Section~\ref{sec:store} invokes this extended $\textsf{processSlot}$ in place of the base one.

\begin{figure}[H]
\caption{Penalty-tracking state extension (leak section)}\label{alg:leak-processslot}
\begin{center}
\begin{boxedminipage}{\textwidth}
\begin{algorithmic}[1]

\Function{$\textsf{processSlot}$}{$\sigma$} \Comment{extended: tracks penalties}
  \State $\sigma.s \gets \sigma.s + 1$
  \State $\sigma_{\mathrm{pre}} \gets \sigma$ \Comment{snapshot of pre-advance state}
  \State $\sigma \gets \textsf{processHeight}(\sigma)$
  \State $\sigma \gets \textsf{processInactivityPenalties}(\sigma_{\mathrm{pre}},\; \sigma)$
  \State \Return $\sigma$
\EndFunction
\vspace{3pt}

\Function{$\textsf{processInactivityPenalties}$}{$\sigma_{\mathrm{pre}},\; \sigma_{\mathrm{post}}$}
  \If{$\sigma_{\mathrm{pre}}.h = \sigma_{\mathrm{post}}.h$} \Comment{Layer 1: stall}
    \For{each $i \in \{1, \ldots, n\}$}
      \If{$\sigma_{\mathrm{post}}.\mathsf{ntargets}[i] = \bot$}
        \State $\sigma_{\mathrm{post}}.\mathsf{penalties}[i] \gets \sigma_{\mathrm{post}}.\mathsf{penalties}[i] + 1$
      \EndIf
    \EndFor
  \EndIf
  \If{$\sigma_{\mathrm{pre}}.J = \sigma_{\mathrm{post}}.J$} \Comment{Layer 2 target: justification did \emph{not} advance}
    \State $T^* \gets \arg\max_{T' \in \{\sigma_{\mathrm{pre}}.\mathsf{jtargets}[i] \,:\, \sigma_{\mathrm{pre}}.\mathsf{jtargets}[i] \neq \bot\}} |\{i : \sigma_{\mathrm{pre}}.\mathsf{jtargets}[i] = T'\}|$ \Comment{$T^* \gets \bot$ if the set is empty; ties broken deterministically}
    \For{each $i \in \{1, \ldots, n\}$}
      \If{$\sigma_{\mathrm{pre}}.\mathsf{jtargets}[i] \neq T^*$ \textbf{or} $\sigma_{\mathrm{pre}}.\mathsf{jtargets}[i] = \bot$}
        \State $\sigma_{\mathrm{post}}.\mathsf{penalties}[i] \gets \sigma_{\mathrm{post}}.\mathsf{penalties}[i] + 1$
      \EndIf
    \EndFor
  \EndIf
  \If{$\sigma_{\mathrm{pre}}.F = \sigma_{\mathrm{post}}.F$ \textbf{and} $\sigma_{\mathrm{pre}}.F \prec \sigma_{\mathrm{pre}}.J$} \Comment{Layer 2 finalize: $F$ did not advance, pending finality}
    \For{each $i \in \{1, \ldots, n\}$}
      \If{$i \notin \sigma_{\mathrm{pre}}.P$}
        \State $\sigma_{\mathrm{post}}.\mathsf{penalties}[i] \gets \sigma_{\mathrm{post}}.\mathsf{penalties}[i] + 1$
      \EndIf
    \EndFor
  \EndIf
  \State \Return $\sigma_{\mathrm{post}}$
\EndFunction

\end{algorithmic}
\end{boxedminipage}
\end{center}
\end{figure}

\paragraph{Reading Figure~\ref{alg:leak-processslot}.}
Figure~\ref{alg:leak-processslot} is the authoritative specification.
$\textsf{processInactivityPenalties}$ checks three independent guards; any subset may fire on a given slot.

\emph{Layer~1 (stall).}
The guard $\sigma_{\mathrm{pre}}.h = \sigma_{\mathrm{post}}.h$ fires when the slot did not advance height.
Every~$i$ with $\sigma_{\mathrm{post}}.\mathsf{ntargets}[i] = \bot$ is bumped, i.e.\ every validator that failed to cast a fresh vote at the current height ($\mathsf{ntargets}$ is not reset at a stall).

\emph{Layer~2 target.}
The guard $\sigma_{\mathrm{pre}}.J = \sigma_{\mathrm{post}}.J$ fires when the justification branch of $\textsf{processHeight}$ did not fire.
$T^*$ is the plurality target across $\sigma_{\mathrm{pre}}.\mathsf{jtargets}$ (deterministic tiebreak), or $\bot$ if no $\just$ votes have been registered.
Every~$i$ with $\sigma_{\mathrm{pre}}.\mathsf{jtargets}[i] \neq T^*$ or $\sigma_{\mathrm{pre}}.\mathsf{jtargets}[i] = \bot$ is bumped (``did not $\just$-vote for the plurality target, or did not $\just$-vote at all'').
When the justification branch fires, it strictly advances~$J$ (Lemma~\ref{lem:slotmono}), so this guard is skipped automatically.
When it does not, the plurality count on $\sigma_{\mathrm{pre}}.\mathsf{jtargets}$ is $< 2n/3$ (counting $T^* = \bot$ as~$0$), so $|\{i : \sigma_{\mathrm{pre}}.\mathsf{jtargets}[i] \neq T^* \vee \sigma_{\mathrm{pre}}.\mathsf{jtargets}[i] = \bot\}| > n/3$; in particular, when $T^* = \bot$ all~$n$ validators are bumped.

\emph{Layer~2 finalize.}
The guard $\sigma_{\mathrm{pre}}.F = \sigma_{\mathrm{post}}.F$ and $\sigma_{\mathrm{pre}}.F \prec \sigma_{\mathrm{pre}}.J$ fires when a pending-finality gap exists at the start of the slot and~$F$ did not advance.
Every $i \notin \sigma_{\mathrm{pre}}.P$ is bumped (``did not commit to finalize the pending justified block'').
When the finality branch of $\textsf{processHeight}$ fires, $\sigma_{\mathrm{post}}.F = \sigma_{\mathrm{pre}}.J \succ \sigma_{\mathrm{pre}}.F$, so this guard is skipped automatically.

\paragraph{Canonical-freshness.}
A vote~$v$ processed at the current slot on the canonical chain~$X$ at state-height $h_c = \sigma_{\mathrm{post}}.h$ bumps $\sigma.\mathsf{ntargets}[v.\validator]$ on~$X$ iff $v$ is \emph{canonical-fresh}:
\[
  v.\height = h_c \;\wedge\; v.\target \preceq X \;\wedge\; v.\target.\slot \geq \sigma_{\mathrm{post}}.s_h,
\]
equivalently $\sigma[v.\target].h = h_c$ with $v.\target \preceq X$ (Remark~\ref{rem:fresh-equiv}).

\subsection{Leak tightness}

\begin{definition}[Non-finality period]\label{def:non-finality-period}
Fix any chain on which slots are processed, and write $\sigma_{\mathrm{pre}}^{(t)}, \sigma_{\mathrm{post}}^{(t)}$ for the pre/post-$\textsf{processHeight}$ snapshots on that chain at slot~$t$.
A \emph{non-finality period} is a maximal interval $[t_1, t_T]$ of consecutive slots processed on the chain such that, at every slot~$t$ in the interval,
\[
  \sigma_{\mathrm{post}}^{(t)}.F = \sigma_{\mathrm{pre}}^{(t)}.F.
\]
Equivalently, $F$ is stuck at some fixed value $F_0$ throughout the interval: the finality branch of $\textsf{processHeight}$ does not fire anywhere in the interval.
Each slot~$t$ in the period falls into exactly one of three categories, determined by its pre-state and by whether $J$ advances:
\begin{itemize}
  \item \emph{pending-stuck}: $\sigma_{\mathrm{pre}}^{(t)}.F \prec \sigma_{\mathrm{pre}}^{(t)}.J$ (a pending-finality gap is open at the start of the slot);
  \item \emph{stable-stuck}: $\sigma_{\mathrm{pre}}^{(t)}.F = \sigma_{\mathrm{pre}}^{(t)}.J$ and $\sigma_{\mathrm{post}}^{(t)}.J = \sigma_{\mathrm{pre}}^{(t)}.J$ (no pending gap, and $J$ does not advance at~$t$);
  \item \emph{transition}: $\sigma_{\mathrm{pre}}^{(t)}.F = \sigma_{\mathrm{pre}}^{(t)}.J$ and $\sigma_{\mathrm{post}}^{(t)}.J \succ \sigma_{\mathrm{pre}}^{(t)}.J$ (no pending gap at the start, but the justification branch fires and opens one by the end of the slot).
\end{itemize}
\end{definition}

\begin{theorem}[Non-finality period penalty bound]\label{thm:tightness}
Over any non-finality period of length $T \geq 1$ on any chain, the total penalty increment across all validators assessed by $\textsf{processInactivityPenalties}$ on that chain is at least $(T - 1) \cdot n/3$.
\end{theorem}

\begin{proof}
Fix a non-finality period $[t_1, t_T]$ with $T \geq 1$.
Every slot $t \in [t_1, t_T]$ is pending-stuck, stable-stuck, or transition by Definition~\ref{def:non-finality-period}.
We show that (i) at most one slot of the period is transition, and (ii) every pending-stuck or stable-stuck slot contributes strictly more than $n/3$ to the total penalty increment.

\emph{(i) At most one transition slot.}
A transition at~$t$ requires $\sigma_{\mathrm{pre}}^{(t)}.F = \sigma_{\mathrm{pre}}^{(t)}.J$ and $\sigma_{\mathrm{post}}^{(t)}.J \succ \sigma_{\mathrm{pre}}^{(t)}.J$.
After such a slot, $\sigma_{\mathrm{post}}^{(t)}.J \succ \sigma_{\mathrm{pre}}^{(t)}.F$; since $F$ does not advance anywhere in the period, $\sigma_{\mathrm{pre}}^{(t')}.F = \sigma_{\mathrm{pre}}^{(t)}.F \prec \sigma_{\mathrm{pre}}^{(t')}.J$ at every later slot $t' > t$ of the period (the justification branch of $\textsf{processHeight}$ only reassigns $J \gets T$ when it strictly advances along~$\preceq$, so~$J$ is monotone across slots).
So every slot after a transition is pending-stuck, and no second transition can occur within the period.

\emph{(ii) Per-slot penalty $> n/3$ for pending-stuck and stable-stuck slots.}

\textbf{Pending-stuck slot~$t$.}
$\sigma_{\mathrm{pre}}^{(t)}.F \prec \sigma_{\mathrm{pre}}^{(t)}.J$ and $\sigma_{\mathrm{post}}^{(t)}.F = \sigma_{\mathrm{pre}}^{(t)}.F$, so the Layer~2 finalize guard in Figure~\ref{alg:leak-processslot} fires at~$t$.
The finality branch of $\textsf{processHeight}$ did not fire at~$t$ (else $\sigma_{\mathrm{post}}^{(t)}.F = \sigma_{\mathrm{pre}}^{(t)}.J \succ \sigma_{\mathrm{pre}}^{(t)}.F$), so $|\sigma_{\mathrm{pre}}^{(t)}.P| < 2n/3$, whence $|\{i : i \notin \sigma_{\mathrm{pre}}^{(t)}.P\}| > n/3$.
Layer~2 finalize charges strictly more than $n/3$ validators at~$t$.

\textbf{Stable-stuck slot~$t$.}
$\sigma_{\mathrm{pre}}^{(t)}.J = \sigma_{\mathrm{post}}^{(t)}.J$, so the Layer~2 target guard in Figure~\ref{alg:leak-processslot} fires at~$t$.
The justification branch of $\textsf{processHeight}$ did not fire at~$t$ (else $J$ would advance), so the plurality count on $\sigma_{\mathrm{pre}}^{(t)}.\mathsf{jtargets}$ is $< 2n/3$ (counting $T^* = \bot$ as~$0$), whence $|\{i : \sigma_{\mathrm{pre}}^{(t)}.\mathsf{jtargets}[i] \neq T^* \vee \sigma_{\mathrm{pre}}^{(t)}.\mathsf{jtargets}[i] = \bot\}| > n/3$.
Layer~2 target charges strictly more than $n/3$ validators at~$t$.

\emph{Conclusion.}
At most one slot of the period is transition, so at least $T - 1$ slots are pending-stuck or stable-stuck.
Each such slot contributes strictly more than $n/3$ by~(ii); Layer~1 and the remaining guards add only non-negative charges.
Summing over the period, the total penalty increment across all validators is at least $(T - 1) \cdot n/3$.

The bound is purely state-logic: the proof reads $\sigma_{\mathrm{pre}}^{(t)}$ and $\sigma_{\mathrm{post}}^{(t)}$ on the fixed chain and counts charges directly from Figure~\ref{alg:leak-processslot}, with no hypothesis on synchrony, on the honest voting rule, or on the size of the slashable set.
\end{proof}

\subsection{Honest voting rule}

We state the honest rule explicitly, as it underpins the leak fairness proof.

\begin{definition}[Honest voting rule]\label{def:voting-rule}
Let $X = \textsf{getHead}(\Sigma)$ be the canonical chain at the voting opportunity, $\sigma[X]$ its post-state, $h_c = \sigma[X].h$ the current canonical height, $J = \sigma[X].J$ the canonical justified block and $h_j = \sigma[X].h_j$ its height.
Write $\mathsf{currentSlot}$ for the node's protocol time at the voting opportunity, a slot counter shared across chains with $\mathsf{currentSlot} = \sigma[X].s \geq \sigma[X].s_h$.
Let $T$ denote the canonical target at height~$h_c$: any ancestor of~$X$ with $\sigma[T].h = h_c$ (e.g.\ $X$ itself).

\emph{Target selection.}
The target $T_c$ of~$i$'s vote at height~$h_c$ is
\begin{itemize}
  \item $T_c = L$ if $i$ is \emph{locked at~$h_c$} --- i.e., has previously cast a vote with $\mathsf{finalize\_height} = h_c$ and $\mathsf{finalize\_target} = L \neq \bot$, forced by E1 since committing~$L$ at~$h_c$ forbids any future height-$h_c$ vote with target $\neq L$;
  \item $T_c = T$ otherwise, read from the current $\sigma[X]$.
\end{itemize}

\emph{Voting rules.}
At each voting opportunity at canonical height~$h_c$, validator~$i$ acts as follows:
\begin{enumerate}
  \item[(R1)] If $i$ has not yet cast any vote at height~$h_c$: cast a $\just$ vote
  \[ v = (i,\; h_c,\; \mathsf{currentSlot},\; T_c,\; \just,\; f_h,\; f_t) \]
  with finalize-fields
  \[
    (f_h,\, f_t) \;=\;
    \begin{cases}
      (h_j,\, J) & \parbox[t]{0.55\textwidth}{\raggedright if every prior vote of~$i$ at height~$h_j$ had target~$J$;}\\[0.8em]
      (\bot,\, \bot) & \text{otherwise.}
    \end{cases}
  \]
  \item[(R2)] Otherwise: cast a $\notar$ vote
  \[ v = (i,\; h_c,\; \mathsf{currentSlot},\; T_c,\; \notar,\; \bot,\; \bot). \]
\end{enumerate}
Both rules evaluate $T$, $J$, $h_j$ against the current $\sigma[X]$; if~$i$'s canonical view shifted between prior firings and the current one, $T_c$ may shift accordingly.
Only (R1) produces a $\just$ vote, and (R1) fires at most once per~$h_c$, so~$i$ casts at most one $\just$ vote per height.
Voting at~$h_c$ stops naturally once canonical advances past~$h_c$.
\end{definition}

Under synchrony, $f < n/3$, and deterministic $\textsf{getHead}$, honest validators share the same canonical state (Corollary~\ref{cor:agreement}) and agree on~$T$, so the (R1) $\just$ quorum for~$T$ forms at~$h_c$ and justification fires in the first slot at~$h_c$.
Rule~(R2) is the \emph{resubmission} fallback for later slots at the same canonical height: when the (R1) $\just$ quorum did not form on canonical (heterogeneous targets across views), the $\notar$ vote's prefix counting may still notarize a common ancestor.
Re-reading canonical at each (R2) ensures the $\notar$ vote tracks~$i$'s current canonical view; the lock clause preserves E1-safety under adversarial views or retries.
Since (R2) fires every slot while canonical stays at~$h_c$, $\mathsf{ntargets}[i]$ remains populated on every chain that keeps processing~$i$'s new votes --- processing old votes is never required for progress, and they can be discarded.

An (R1) vote optionally commits to finalize~$J$ at~$h_j$ via $(f_h, f_t) = (h_j, J)$, placing~$i$ in~$P$ once $(J, h_j)$ is processed.
The finalize-field conditional checks \emph{all} prior votes of~$i$ at height~$h_j$, not just the (R1) one: a (R2) at~$h_j$ may have chosen a target differing from the (R1) target if~$i$'s canonical view shifted between firings.
Committing $f_t = J$ at height~$h_j$ requires (by E1) that every vote of~$i$ at~$h_j$ have target~$J$; if (R1) and (R2) at~$h_j$ chose different targets, the conditional fires $(\bot, \bot)$ and~$i$ misses the finalization at~$h_j$ but remains E1-safe.

\paragraph{E1 safety of the rule.}
The rule is E1-safe: Lemma~\ref{lem:honest-e1-safety} records that an honest validator never triggers E1 slashing.
Per-height $\just$-uniqueness is immediate: only (R1) produces a $\just$ vote, and (R1) fires at most once per~$h_c$.

\begin{lemma}[Honest E1 safety]\label{lem:honest-e1-safety}
Under synchrony, $f < n/3$, and all honest validators following Definition~\ref{def:voting-rule}, an honest~$i$ never produces a vote-pair $(v_1, v_2)$ (possibly $v_1 = v_2$) with $v_2.\mathsf{finalize\_target} \neq \bot$, $v_1.\height = v_2.\mathsf{finalize\_height}$, and $v_1.\target \neq v_2.\mathsf{finalize\_target}$.
\end{lemma}

\begin{proof}
Only (R1) produces a finalize-bearing vote, with $(f_h, f_t) = (h_j, J)$ read from the canonical state at~$v_2$'s firing.
Write $h_* = v_2.\mathsf{finalize\_height} = h_j$ and $T_* = v_2.\mathsf{finalize\_target} = J$.
We show every vote~$v_1$ of~$i$ with $v_1.\height = h_*$ has $v_1.\target = T_*$, by cases on the temporal relation between $v_1$ and $v_2$.

\emph{Case~A ($v_1$ cast before~$v_2$).}
The (R1) finalize-field conditional at~$v_2$ attached $(h_*, T_*)$ only because every prior vote of~$i$ at height~$h_*$ had target~$T_*$; in particular, $v_1.\target = T_*$.

\emph{Case~B ($v_1$ cast after~$v_2$).}
If $h_* = 0$ (so $v_2$ attached its finalize commitment before any justification had fired, with $T_* = \genesis$), the lock clause of Definition~\ref{def:voting-rule} pins $T_c = \genesis$ at any later height-$0$ voting opportunity of~$i$, so $v_1.\target = \genesis = T_*$; the case reduces to target-matching.
Otherwise $h_* \geq 1$, so $v_2$ fired at canonical height $h_c > h_*$.
Under the hypotheses, every canonical height-advance is witnessed by the $> 2n/3$ honest (R1) quorum agreeing on~$T$ (Corollary~\ref{cor:agreement}), so canonical state-height is monotone over time; canonical cannot return to height~$h_*$ for any later~$v_1$, and the case is vacuous.

\emph{Case~C ($v_1 = v_2$).}
Vacuous: $v_2.\height = h_c \neq h_* = v_2.\mathsf{finalize\_height}$.
\end{proof}

\begin{remark}[Scope of Lemma~\ref{lem:honest-e1-safety}]\label{rem:honest-e1-scope}
All three hypotheses are load-bearing for Case~B.
A validator that wishes not to rely on them may always set $(f_h, f_t) = (\bot, \bot)$ at (R1), trivially avoiding E1 at the cost of not contributing to finalization.
\end{remark}

\subsection{Leak fairness}

Fairness is local to a single honest validator~$i$: we fix~$i$, its local store~$\Sigma$, and the canonical chain $X = \textsf{getHead}(\Sigma)$ that~$i$ sees.
The claim is that the vote~$v$ produced at the current slot by the honest rule (Definition~\ref{def:voting-rule}), once included on any extension of~$X$ still at the same state-height, exempts~$i$ from Layer~1 penalties at that state-height on that extension.
Layer-2 exemption is out of scope: matching the plurality target when justification does not advance, and entering~$P$ before pending finality closes, both depend on timely quorum formation rather than on any local property of~$i$.
Because height advances reset $\mathsf{jtargets}$ and $\mathsf{ntargets}$ (Figure~\ref{alg:state-machine}), the Layer-2 state at slot~$s+1$ is independent of~$i$'s votes cast at prior heights, so no ``lock'' issue obstructs fairness at the next height.

\begin{theorem}[Layer~1 leak fairness]\label{thm:fairness}
Unless $\geq n/3$ validators are slashable, and assuming all honest validators follow Definition~\ref{def:voting-rule}: let~$i$ be an honest validator, $\Sigma$ its local store, $X = \textsf{getHead}(\Sigma)$ its canonical chain, and $h_c = \Sigma.\sigma[X].h$ the current canonical state-height.
Let $v$ be the vote~$i$ produces at the current slot.
Let $\Sigma'$ be any store obtained by processing additional blocks and votes on top of~$\Sigma$ such that $X' = \textsf{getHead}(\Sigma')$ is an extension of~$X$ with $\Sigma'.\sigma[X'].h = h_c$, and suppose~$v$ is included on~$X'$.
Then for every slot processed on~$X'$ from that point onward while canonical state-height stays at~$h_c$, the Layer~1 increment to $\sigma.\mathsf{penalties}[i]$ on canonical is zero.
\end{theorem}

\begin{proof}
It suffices to show that the target $T_c$ selected by Definition~\ref{def:voting-rule} at~$i$'s voting opportunity satisfies $T_c \preceq X$ with $\sigma[T_c].h = h_c$: freshness (Remark~\ref{rem:fresh-equiv}) then implies~$v$ is canonical-fresh on~$X'$, processing~$v$ populates $\mathsf{ntargets}[i]$ on~$X'$, and Layer~1 does not bump~$i$ at any subsequent slot while state-height stays at~$h_c$.

We identify the top-tier block~$N$ of the walk that produced $\Sigma.R_s$, bound $h_c$ in terms of its height~$h$, and handle the locked branch of $T_c$ (if reached) by matching its target against~$N$.

\emph{Step 1: top-tier block~$N$.}
Let $h \geq 0$ be the maximum height across all tuples in $\Sigma.C_N$ (well-defined since $(\genesis, 0, 0) \in C_N$), and let~$N$ be the unique block such that $(N, h, s_N) \in \Sigma.C_N$ is lex-maximum over $\{(X', h, s) \in C_N\}$ under the key $(h, \mathbf{1}[\cdot \in C_J], s, \mathsf{hash}(\cdot))$ (uniqueness: block hashes are unique).
If any height-$h$ tuple lies in~$C_J$, the indicator bit forces~$N$'s tuple into~$C_J$.

\emph{Step 2: $N \preceq R_s$.}
If $h = 0$, then $N = \genesis \preceq R_s$ by Theorem~\ref{thm:fleqr}.
Otherwise $h \geq 1$.
Let $h_{F_s}$ be the height at which~$F_s$ was finalized (or~$0$ if $F_s = \genesis$).
Remark~\ref{rem:fs-invariant} gives $(F_s, h_{F_s}, F_s.\slot) \in C_J \subseteq C_N$, so maximality of~$h$ yields $h \geq h_{F_s}$; Lemma~\ref{lem:certchain} then gives $F_s \sim N$.
Split on this compatibility ($F_s = N$ falls into the first branch by reflexivity of~$\preceq$):
\begin{itemize}
  \item $N \preceq F_s$: Theorem~\ref{thm:fleqr} gives $N \preceq F_s \preceq R_s$.
  \item $F_s \prec N$: $N$'s tuple is a strict-descendant candidate at the first step of $\textsf{computeRoot}$ from~$F_s$.
  Any other strict-descendant candidate $(X', h', s') \in C_N$ has $h' \leq h$ by maximality; if $h' < h$ its key is lex-dominated by~$N$'s, and if $h' = h$ Step~1 picks~$N$'s tuple as lex-max.
  Either way, the walk's first step picks~$N$'s tuple, so $N \preceq R_s$.
\end{itemize}

\emph{Step 3: $N \preceq X$ and $h_c \geq h$.}
$X$ descends from~$R_s$ (Figure~\ref{alg:store}), so $R_s \preceq X$ and thus $N \preceq X$.
Remark~\ref{rem:fresh-equiv} gives $\sigma[N].h = h$, and state-height monotonicity along~$\preceq$ (Lemma~\ref{lem:slotmono}) gives $h_c = \sigma[X].h \geq h$.

\emph{Step 4: $h_c \leq h + 1$.}
Every state-height advance on~$X$ from some $h'$ to~$h' + 1$ is gated by a height-$h'$ quorum (Lemma~\ref{lem:heightprog}), whose processing inserts a tuple at height~$h'$ into~$C_N$.
Since~$h$ is the maximum height in~$C_N$, no advance from $h + 1$ to~$h + 2$ has fired on~$X$.
With Step~3, $h_c \in \{h,\, h+1\}$.

\emph{Characterization of lock at height~$h'$.}
Before splitting on~$h_c$, we record how locks arise.
A lock at height~$h'$ requires a prior (R1) firing of~$i$ with finalize-fields $(h', J_\ast)$, where $J_\ast = \sigma[X_\ast].J$ and $h' = \sigma[X_\ast].h_j$ for the canonical chain $X_\ast$ at that firing.
At that firing, the ancestor of~$X_\ast$ whose processing last advanced~$J$ to~$J_\ast$ inserted $(J_\ast, h', J_\ast.\slot)$ into~$C_J$, and $C_J$ is only inserted into (Figure~\ref{alg:store}).
So a lock at height~$h'$ entails a height-$h'$ tuple in $\Sigma.C_J$, with block exactly~$J_\ast$.

\emph{Step 5: case $h_c = h + 1$.}
By the maximality of~$h$, $\Sigma.C_J \subseteq \Sigma.C_N$ has no height-$(h+1)$ tuple, so by the characterization no lock at~$h_c$ exists.
Hence $T_c = T \preceq X$ with $\sigma[T].h = h_c$ by Remark~\ref{rem:fresh-equiv}.

\emph{Step 6: case $h_c = h$, with a height-$h$ $C_J$-tuple in~$\Sigma.C_J$.}
By Step~1's tiebreak,~$N$'s tuple lies in~$C_J$, so~$N$ was justified at height~$h$ on some chain.

\emph{Uniqueness of height-$h$ justified blocks.}
Every honest validator casts at most one $\just$ vote at any height (Definition~\ref{def:voting-rule}).
Let $Q_1, Q_2$ be $\just$-quora at height~$h$ for blocks $N_1, N_2$ respectively.
$|Q_1 \cap Q_2| \geq 2n/3 + 2n/3 - n = n/3 > f$, so the intersection contains an honest validator~$j$, who cast $\just$ votes targeting both $N_1$ (contributing to $Q_1$) and $N_2$ (contributing to $Q_2$); honesty forces these to be the same vote, hence $N_1 = N_2$.
Every height-$h$ $C_J$-tuple therefore has block~$N$.

Applying the characterization: if~$i$ is locked at~$h_c = h$, then $T_c = L$ where~$L$ is the $J_\ast$ of the locking (R1), and uniqueness forces $L = N$; with Step~3, $T_c = N \preceq X$ and $\sigma[T_c].h = h = h_c$.
If~$i$ is not locked, $T_c = T \preceq X$ with $\sigma[T].h = h_c$ directly.

\emph{Step 7: case $h_c = h$, with no height-$h$ $C_J$-tuple in~$\Sigma.C_J$.}
By the characterization, no lock at~$h_c = h$ can have been created, so~$i$ is not locked.
Then $T_c = T \preceq X$ with $\sigma[T].h = h_c$.

\emph{Step 8: canonical-freshness and Layer~1.}
Each case gives $T_c \preceq X$ with $\sigma[T_c].h = h_c$.
Since $X \preceq X'$, also $T_c \preceq X'$.
The hypothesis $\Sigma'.\sigma[X'].h = h_c = \Sigma.\sigma[X].h$ and Lemma~\ref{lem:slotmono} ($s_h$ advances only with~$h$) give $\Sigma'.\sigma[X'].s_h = \Sigma.\sigma[X].s_h$, and Remark~\ref{rem:fresh-equiv} gives $T_c.\slot \geq \Sigma'.\sigma[X'].s_h$.
The produced vote $v = (i, h_c, \mathsf{currentSlot}, T_c, \cdot, \ldots)$ is therefore canonical-fresh on~$X'$, and processing~$v$ sets $\Sigma'.\sigma[X'].\mathsf{ntargets}[i] \gets \max_{\preceq}(\Sigma'.\sigma[X'].\mathsf{ntargets}[i],\, T_c) \neq \bot$.

$\mathsf{ntargets}$ is reset only at a height advance (Figure~\ref{alg:state-machine}), and no such advance fires on~$X'$ while canonical state-height stays at~$h_c$, so $\mathsf{ntargets}[i] \neq \bot$ persists on canonical from the inclusion slot onward.
Layer~1 bumps only validators with $\sigma_{\mathrm{post}}.\mathsf{ntargets}[\cdot] = \bot$, so it does not bump~$i$.
\end{proof}

The hypothesis is used in Step~6 (uniqueness of height-$h$ justified blocks) and implicitly in Step~2 (Lemma~\ref{lem:certchain}); together these force the locked-branch target~$L$ to coincide with the walk's~$N$, hence to lie on canonical.

\section{Roberto's review}\label{sec:roberto-review}

This section collects the revisions and additions proposed during review.
The main text above is preserved as written; each entry below either supersedes a corresponding item in the main text (indicated via back-references) or introduces new material.
Labels here are prefixed \texttt{rev:} to avoid clashing with the main-text labels.

\subsection{Store root via a single argmax}\label{rev:sec:store-root}

\paragraph{Motivation.}
The revision to $\textsf{computeRoot}$ is motivated by the development in \Cref{rev:sec:leak-fairness} and \Cref{sec:healing}, specifically \Cref{rev:lem:root-height-cj-exclusive} (``root-height exclusivity''), which states that if $R_s \in C_N \setminus C_J$ then no $(B, h, s) \in C_J$ has $h = R_s.h$.
This lemma fails under the iterative walk: if $(A, h, \cdot) \in C_J$ and $(B, h, \cdot) \in C_N \setminus C_J$ with $A \prec B$ at the same state-height, the walk picks~$A$ at its first step (bit tiebreak), then advances to~$B$ as a strict descendant, producing $R_s = B \in C_N \setminus C_J$ while $A \in C_J$ sits at the same height --- witnessing the negation.
Replacing the walk by a single argmax resolves this: the argmax stops at~$A$ (the $C_J$-bit breaks the tie in $A$'s favour), so $R_s = A \in C_J$ and the ``$R_s \in C_N \setminus C_J$'' hypothesis of \Cref{rev:lem:root-height-cj-exclusive} is vacuous against a same-height $C_J$ competitor.
\Cref{rev:lem:root-height-cj-exclusive} underpins \Cref{rev:lem:canonical-fresh-honest-vote} and hence the healing lemmas; the closely related \Cref{rev:lem:Rs-start-anywhere} and \Cref{rev:cor:no-high-cn} likewise rely on $R_s$ being the argmax itself rather than a descendant.

\paragraph{Revised \texttt{computeRoot}.}
The iterative greedy walk in Figure~\ref{alg:store} is replaced by a single argmax:
\begin{center}
\begin{minipage}{0.9\textwidth}
\begin{algorithmic}[1]
\Function{$\textsf{computeRoot}$}{$\Sigma$} \Comment{Single argmax from $\Sigma.F_s$}
  \State $(X^*, \_, \_) \gets \arg\max_{(X,h,s) \in \Sigma.C_N,\; X \succeq \Sigma.F_s} \bigl(h,\; \mathbf{1}[(X,h,s) \in \Sigma.C_J],\; s,\; \mathsf{hash}(X)\bigr)$
  \State \Return $X^*$
\EndFunction
\end{algorithmic}
\end{minipage}
\end{center}

\paragraph{Revised proofs of downstream theorems.}
The following proofs are revised to use the single-argmax form. Only the proofs change; the statements of \Cref{thm:leafequiv,thm:fleqr,lem:upgrade,thm:orderindep} are unchanged.

\emph{Revised proof of \Cref{thm:leafequiv}.}
Let $\tau^\star = (X^\star, h^\star, X^\star.\slot)$ be the full argmax over $\{(X, h, s) \in C_N : X \succeq F_s\}$, so $\textsf{computeRoot}$ returns $X^\star$.
We show $\tau^\star \in C_N^{\mathrm{leaf}}$, which implies $\textsf{computeRootLeaf}$ picks the same tuple.
Let $B^\star$ be any processed block whose post-state contributed~$\tau^\star$; if $B^\star$ is a leaf, done.
Otherwise pick any processed leaf $B_\ell \succ B^\star$; its post-state contributes $\tau_\ell = (X_\ell, h_\ell, X_\ell.\slot) \in C_N^{\mathrm{leaf}}$.
By \Cref{lem:slotmono} along the chain $B^\star$ to $B_\ell$, $h_\ell \geq h^\star$ and $X_\ell \succeq X^\star$, so $\tau_\ell$ is a candidate.
If $h_\ell > h^\star$, $\tau_\ell$'s key strictly dominates $\tau^\star$'s by the leading coordinate, contradicting maximality.
Hence $h_\ell = h^\star$, and the invariant that $N$ advances only jointly with $h_n \gets h_n + 1$ forces $X_\ell = X^\star$, so $\tau_\ell = \tau^\star \in C_N^{\mathrm{leaf}}$.

\emph{Revised proof of \Cref{thm:fleqr}.}
Initially $F_s = R_s = \genesis$.
$\textsf{computeRoot}$ returns the block of a tuple with $X \succeq F_s$, so $R_s \succeq F_s$ holds just after $\textsf{computeRoot}$.
$\textsf{updateFinalized}$ sets $F_s$ only to values $F_B \preceq R_s$, preserving the invariant.

\emph{Revised proof of \Cref{lem:upgrade}.}
Since $(F, h_F, F.\slot) \in C_J \subseteq C_N$ and $F \succ F_s$, it is a candidate for $\textsf{computeRoot}$'s argmax, so the returned tuple has key $\geq (h_F, 1, F.\slot, \mathsf{hash}(F))$ and $h \geq h_F$.
If $h > h_F$, \Cref{lem:certchain} gives $F \prec X$.
If $h = h_F$, the key lower bound forces the $C_J$-bit to be~$1$, i.e.\ $(X, h, s) \in C_J$; the slot/hash argument of the original proof gives $F \preceq X$.
Either way $R_s = X \succeq F$.

\emph{Revised proof of \Cref{thm:orderindep}.}
Same argument as the original proof for $\sigma$, $C_J$, $C_N$, and $F_s$; the final step becomes: $R_s$ is the argmax over $\{(X,h,s) \in C_N : X \succeq F_s\}$ on the totally ordered key, hence a deterministic function of $(C_N, C_J, F_s)$.

\subsection{Unified target selection in the honest voting rule}\label{rev:sec:voting-rule}

\paragraph{Status: no longer required.}
An earlier version of this review proposed unifying (R1) and (R2) to share a single target~$T_c$, defined once before the case split:
\[
  T_c \;=\;
  \begin{cases}
    v'.\mathsf{finalize\_target} & \text{if $i$ has previously cast a vote $v'$ with $v'.\mathsf{finalize\_height} = h_c$;}\\
    T & \text{otherwise (the canonical target).}
  \end{cases}
\]
This unification was motivated by the counterexample to honest E1 safety under asynchrony: without the lock on~(R1)'s target, a validator whose canonical view fork-switched back to height~$h_c$ could cast a~(R1) $\just$ for a fresh target~$T$ while still holding an earlier $\mathsf{finalize\_target} \neq T$ at~$h_c$, triggering E1.

\Cref{def:voting-rule} in the current main text already incorporates an equivalent fix via its ``locked at~$h_c$'' clause, which applies the lock to \emph{both} (R1) and (R2).
The proposal in this review is therefore redundant with the main text and is retained here only to record the motivation and the \Cref{rev:sec:async-e1} lemma that it enables.

\subsection{Asynchronous honest E1 safety}\label{rev:sec:async-e1}

The ``locked at $h_c$'' clause in \Cref{def:voting-rule} --- ``$T_c = L$ whenever $i$ has previously cast a vote with $\mathsf{finalize\_height} = h_c$ and $\mathsf{finalize\_target} = L \neq \bot$'' --- makes the asynchronous version of \Cref{lem:honest-e1-safety} provable without the synchrony hypothesis that Case~B of the original proof relied on.

\begin{lemma}[Honest E1 safety, asynchronous]\label{rev:lem:honest-e1-safety-asynchrony}
Under $f < n/3$ and all honest validators following \Cref{def:voting-rule}, an honest~$i$ never produces a vote-pair $(v_1, v_2)$ (possibly $v_1 = v_2$) with $v_2.\mathsf{finalize\_target} \neq \bot$, $v_1.\height = v_2.\mathsf{finalize\_height}$, and $v_1.\target \neq v_2.\mathsf{finalize\_target}$.
\end{lemma}

\begin{proof}
Write $h_\ast \coloneqq v_2.\mathsf{finalize\_height}$ and $T_\ast \coloneqq v_2.\mathsf{finalize\_target} \neq \bot$.
Only (R1) produces finalize-bearing votes, so $v_2$ is an (R1) vote with $(f_h, f_t) = (h_j, J) = (h_\ast, T_\ast)$ read from the canonical state at $v_2$'s firing.

\emph{Case~A ($v_1$ cast before $v_2$).}
The (R1) finalize-field conditional at~$v_2$ attached $(h_\ast, T_\ast)$ only because every prior vote of~$i$ at height~$h_\ast$ had target~$T_\ast$; in particular $v_1.\target = T_\ast$.

\emph{Case~B ($v_1$ cast after $v_2$).}
At $v_1$'s firing, the voting height is $h_c = v_1.\height = h_\ast$, and~$i$ has the prior vote~$v_2$ with $\mathsf{finalize\_height} = h_\ast = h_c$ and $\mathsf{finalize\_target} = T_\ast \neq \bot$: $i$ is \emph{locked at $h_c$}, and the target-selection clause of \Cref{def:voting-rule} forces $T_c = T_\ast$ regardless of~$i$'s canonical view at that moment.
Hence $v_1.\target = T_c = T_\ast$.

\emph{Case~C ($v_1 = v_2$).}
Vacuous: $v_2.\height = h_c \neq h_\ast = v_2.\mathsf{finalize\_height}$ (every height advance sets $h_j \gets h$ before $h \gets h + 1$, so $h_j < h$ in any post-advance state).
\end{proof}

The proof uses neither synchrony, nor determinism of $\textsf{getHead}$, nor any agreement on the canonical view across validators --- the lock clause carries Case~B on its own.

\subsection{Additional store-level lemmas}\label{rev:sec:store-lemmas}

\begin{lemma}[Monotonicity of the root height]\label{rev:lem:Rs-height-monotone}
$R_s.h$ is non-decreasing over time.
\end{lemma}

\begin{proof}
$R_s$ changes only inside $\textsf{onBlock}$, which (i) inserts into $C_N, C_J$, (ii) recomputes $R_s$, (iii) possibly advances $F_s$ under the guard $F_s \preceq R_s$.
Between two calls, $F_s^{\text{new}} \preceq R_s^{\text{old}}$ (either unchanged or just-updated under guard), so the old $R_s$-tuple remains a valid argmax candidate at the next $\textsf{computeRoot}$, giving $R_s^{\text{new}}.h \geq R_s^{\text{old}}.h$.
\end{proof}

\begin{lemma}[Root independent of the argmax starting point]\label{rev:lem:Rs-start-anywhere}
Unless $\geq n/3$ validators are slashable, the block~$R_s$ returned by $\textsf{computeRoot}$ coincides with the block of the argmax over $\{(X,h,s) \in C_N : X \succeq \genesis\}$.
\end{lemma}

\begin{proof}
Let $R_s^{\genesis}$ be that block; $\{X \succeq F_s\} \subseteq \{X \succeq \genesis\}$ gives $R_s^{\genesis}.h \geq R_s.h \geq h_{F_s}$.
If $R_s^{\genesis}.h > h_{F_s}$, \Cref{lem:certchain} gives $F_s \prec R_s^{\genesis}$, so $R_s^{\genesis} \succeq F_s$ and $R_s^{\genesis} = R_s$ by uniqueness.
If $R_s^{\genesis}.h = h_{F_s}$, the $(F_s, h_{F_s}, \cdot) \in C_J$ candidate forces the $C_J$-bit of $R_s^{\genesis}$ to be~$1$; \Cref{lem:finchain}-style quorum intersection with $F_s$ finalized gives $R_s^{\genesis} = F_s$, hence again $R_s^{\genesis} = R_s$.
\end{proof}

\begin{corollary}[No $C_N$ record above the root]\label{rev:cor:no-high-cn}
Unless $\geq n/3$ validators are slashable, every $(X, h, s) \in C_N$ satisfies $h \leq R_s.h$; in particular the same holds for every $(X, h, s) \in C_J \subseteq C_N$.
\end{corollary}

\begin{proof}
By \Cref{rev:lem:Rs-start-anywhere}, $R_s$'s tuple is the argmax over all of $C_N$; a $C_N$ tuple at height $> R_s.h$ would strictly dominate on the leading coordinate, contradicting maximality.
\end{proof}

\begin{corollary}[Processed-block height bound]\label{rev:cor:processed-height-bound}
Unless $\geq n/3$ validators are slashable, every processed block~$B$ satisfies $\sigma[B].h \leq R_s.h + 1$.
\end{corollary}

\begin{proof}
If $\sigma[B].h \geq 1$, \Cref{lem:heightprog} produces a $C_N$ tuple at height $\sigma[B].h - 1$ from the last advance on~$B$'s chain; by \Cref{rev:cor:no-high-cn}, $\sigma[B].h - 1 \leq R_s.h$.
\end{proof}

\subsection{Leak fairness 2}\label{rev:sec:leak-fairness}

\begin{lemma}[Unique justified block per height]\label{rev:lem:unique-just-per-height}
If $f < n/3$ and all honest validators follow \Cref{def:voting-rule}, then for any two blocks $C_1, C_2$ justified at the same height~$h$, $C_1 = C_2$.
\end{lemma}

\begin{proof}
Each justification yields a $\just$-quorum $Q_k$ ($k \in \{1,2\}$) of size $\geq 2n/3$ signing a $\just$ vote at height~$h$ for $C_k$; quorum intersection gives $|Q_1 \cap Q_2| \geq n/3$.
Since $f < n/3$, some $v \in Q_1 \cap Q_2$ is honest; by \Cref{def:voting-rule}, only (R1) emits a $\just$ vote and it fires at most once per height, so~$v$'s single $\just$ vote at~$h$ has both $C_1$ and $C_2$ as target, forcing $C_1 = C_2$.
\end{proof}

\begin{lemma}[Root-height exclusivity]\label{rev:lem:root-height-cj-exclusive}
Unless $\geq n/3$ slashable: if $(R_s, R_s.h, \cdot) \in C_N \setminus C_J$ and $(B, h, s) \in C_J$, then $h \neq R_s.h$.
\end{lemma}

\begin{proof}
Write $h_R \coloneqq R_s.h$. Since $(F_s, h_{F_s}, F_s.\slot) \in C_J$ is always an argmax candidate, $R_s \neq F_s$, so $R_s \succ F_s$ and $h_R \geq h_{F_s}$.
Suppose $(B, h_R, s) \in C_J$ exists.
If $h_R > h_{F_s}$: \Cref{lem:certchain} gives $F_s \prec B$, so $B$ is a candidate whose key $(h_R, 1, \ldots)$ beats $R_s$'s $(h_R, 0, \ldots)$ on the $C_J$-bit, contradicting maximality.
If $h_R = h_{F_s}$: $B$ is justified at $h_{F_s}$ and $F_s$ is finalized at $h_{F_s}$, so $B = F_s$ by \Cref{lem:finchain}; then $(F_s, h_{F_s}, F_s.\slot)$ is itself an argmax candidate and wins on the $C_J$-bit, again contradicting maximality.
\end{proof}

\begin{lemma}[No finalize vote without a $C_J$ record]\label{rev:lem:no-finalize-without-cj}
Let~$i$ be a validator with local store~$\Sigma$ and $h$ a height with $(X, h, s) \notin \Sigma.C_J$ for all $X, s$. Then~$i$ has not cast a vote~$v$ with $v.\mathsf{finalize\_height} = h$.
\end{lemma}

\begin{proof}
Only (R1) produces $\mathsf{finalize\_height} \neq \bot$, so such a~$v$ came from (R1) at an earlier moment with store $\Sigma_0$, yielding $\sigma[X_0].h_j = h$ on~$i$'s then-canonical $X_0 = \textsf{getHead}(\Sigma_0)$.
A state has $h_j = h$ only after justification at~$h$ fired on its chain, so the block~$B \preceq X_0$ where it fired has $\sigma[B].J$ at height~$h$ and $\textsf{onBlock}(B)$ inserts $(\sigma[B].J, h, \cdot) \in \Sigma_0.C_J$; monotonicity of $C_J$ forces this tuple into $\Sigma.C_J$, contradicting the hypothesis.
\end{proof}

\begin{lemma}[Head state-height range]\label{rev:lem:head-height-range}
Let $(R_s, R_s.h, \cdot) \in C_N$. Then $\sigma[\textsf{getHead}(\Sigma)].h \in \{R_s.h,\, R_s.h + 1\}$.
\end{lemma}

\begin{proof}
Write $h_R \coloneqq R_s.h$ and let $Y = \textsf{getHead}(\Sigma)$, so $Y \succeq R_s$.
State-height is non-decreasing along $\preceq$ and $\sigma[R_s].h = h_R$, so $\sigma[Y].h \geq h_R$.
If $\sigma[Y].h \geq h_R + 2$, \Cref{lem:heightprog} produces a second height advance on the chain from $R_s$ to~$Y$, inserting a $C_N$ tuple at height $h_R + 1$ descending from $R_s$; its key beats $R_s$'s on the leading coordinate, contradicting $R_s$'s maximality.
\end{proof}

\begin{lemma}[Canonical-fresh honest vote]\label{rev:lem:canonical-fresh-honest-vote}
Assume $f < n/3$ and all honest validators follow \Cref{def:voting-rule}.
Let~$i$ be an honest validator, $\Sigma$ its local store, $H = \textsf{getHead}(\Sigma)$, and $h_c = \sigma[H].h$.
Let~$v$ be the vote~$i$ produces at the current slot.
Then~$v$ is canonical-fresh on every extension~$H'$ of~$H$ satisfying $\sigma[H'].h = h_c$.
\end{lemma}

\begin{proof}
By \Cref{def:voting-rule}, $v.\height = h_c$ and $v.\target = T_c$, where $T_c$ is either the canonical target $T \preceq H$ or $L \neq \bot$ from a prior locked vote at~$h_c$.
It suffices to show $T_c \preceq H$; split on $h_c$ versus $R_s.h$.

\emph{Case $h_c > R_s.h$.}
Any $(Z, h_c, s) \in C_N$ would, via \Cref{lem:certchain}, satisfy $Z \succeq F_s$, making it an argmax candidate with key $(h_c, \ldots)$ beating $R_s$'s $(R_s.h, \ldots)$ --- contradiction.
Hence no $C_N$ tuple at $h_c$ exists, a fortiori none in $C_J$; by \Cref{rev:lem:no-finalize-without-cj} no prior~$i$-vote has $\mathsf{finalize\_height} = h_c$, so the lock clause does not fire and $T_c = T \preceq H$.

\emph{Case $h_c = R_s.h$ and $(R_s, h_c, \cdot) \notin C_J$.}
\Cref{rev:lem:root-height-cj-exclusive} rules out any $C_J$ tuple at~$h_c$; \Cref{rev:lem:no-finalize-without-cj} then gives $T_c = T \preceq H$.

\emph{Case $h_c = R_s.h$ and $(R_s, h_c, \cdot) \in C_J$.}
If the lock clause does not fire, $T_c = T \preceq H$.
Otherwise~$i$ is locked at~$h_c$ with lock target $L$, and some prior vote of~$i$ had $\mathsf{finalize\_target} = L$ justified at~$h_c$; \Cref{rev:lem:unique-just-per-height} and the case hypothesis give $L = R_s$, and $R_s \preceq H$ (\Cref{thm:fleqr}), so $T_c = L = R_s \preceq H$.
\end{proof}

\begin{remark}[Weakening $f < n/3$ under a double-justify slashing rule]\label{rev:rem:canonical-fresh-slashable}
The $f < n/3$ hypothesis in \Cref{rev:lem:canonical-fresh-honest-vote} is needed only by the invocation of \Cref{rev:lem:unique-just-per-height}, whose proof uses \Cref{def:voting-rule}'s behavioural constraint (at most one $\just$ vote per height) to close quorum intersection.
If \Cref{def:slashing} is extended with an additional slashing rule declaring any pair of $\just$ votes at the same height with different targets slashable, then the intersection of two $\just$-quorums at height~$h$ with different targets forces every shared signer to be slashable; assuming fewer than $n/3$ validators are slashable yields one non-slashable signer, whose single $\just$ vote fixes $C_1 = C_2$.
Under this strengthened slashing definition, $f < n/3$ in the hypotheses of \Cref{rev:lem:unique-just-per-height,rev:lem:canonical-fresh-honest-vote} (and downstream lemmas that inherit it) may be replaced by the weaker ``unless $\geq n/3$ validators are slashable.''
\end{remark}

\section{Healing from asynchrony}\label{sec:healing}

This section builds on the revisions in \Cref{sec:roberto-review}.
Throughout, \Cref{def:voting-rule} and \Cref{alg:store} refer respectively to the original definition and to the algorithm figure with the revised $\textsf{computeRoot}$ of \Cref{rev:sec:store-root}, and the auxiliary lemmas of \Cref{rev:sec:store-lemmas,rev:sec:leak-fairness} are used without further comment.

To state the conditions under which the protocol \emph{heals} after a prolonged period of asynchrony, we enrich the system model with slot-based timing and a per-slot view-merge mechanism.
The view-merge mechanism is detailed in \Cref{app:view-merge}; in this section we abstract it via an assumption sufficient for the healing arguments.

\begin{definition}[Slots and voting time]\label{rev:def:slots}
Physical time is partitioned into an infinite sequence of discrete \emph{slots} indexed by $t \in \mathbb{N}$.
Slot~$t$ has a \emph{start time} $\mathsf{start}(t) \in \mathbb{R}_{\geq 0}$ with $\mathsf{start}(t+1) \geq \mathsf{start}(t) + 2\Delta$, and a \emph{voting time} $\mathsf{vote}(t) \coloneqq \mathsf{start}(t) + \Delta$ at which every honest validator~$i$ invokes \Cref{def:voting-rule} on its local store~$\Sigma_i$ to produce at most one slot-$t$ vote.
\end{definition}

\begin{definition}[Syncer]\label{rev:def:syncer}
A publicly known function $\mathsf{syncer} \colon \mathbb{N} \to \{1, \ldots, n\}$ assigns to each slot~$t$ a validator $\mathsf{syncer}(t)$.
During the portion of slot~$t$ preceding $\mathsf{vote}(t)$, $\mathsf{syncer}(t)$ runs the \emph{view-merge protocol} (\Cref{app:view-merge}).
\end{definition}

\begin{assumption}[View-merge guarantee]\label{rev:ass:view-merge}
If $\mathsf{syncer}(t)$ is honest, then there exists a store $\Sigma^\star_t$ such that every honest validator~$i$ has $\Sigma_i = \Sigma^\star_t$ at time $\mathsf{vote}(t)$.
\end{assumption}

By \Cref{thm:orderindep} applied to $\Sigma^\star_t$, if additionally $f < n/3$ and $\textsf{getHead}$ is deterministic, the components $(R_s, F_s, C_J, C_N)$ and the canonical chain $H = \textsf{getHead}(\Sigma^\star_t)$ are common to every honest validator at $\mathsf{vote}(t)$.

\begin{lemma}[Healing: one-slot advance]\label{rev:lem:healing-one-slot}
Assume:
\begin{enumerate}
  \item $f < n/3$ and all honest validators follow \Cref{def:voting-rule};
  \item $\textsf{getHead}$ is deterministic: equal stores yield equal canonical heads;
  \item $\mathsf{start}(t) - \Delta \geq \GST$;
  \item $\mathsf{syncer}(t)$ is honest; let $\Sigma^\star_t$ be the common store, $H \coloneqq \textsf{getHead}(\Sigma^\star_t)$, and $h^\star \coloneqq R_s^{(t)}.h$;
  \item $H.h = h^\star + 1$;
  \item no honest validator has cast any vote at height $> h^\star$ at any time strictly before $\mathsf{vote}(t)$;
  \item by $\mathsf{start}(t+1) - \Delta$, in every honest validator's view, there exists a chain $X \succeq H$ such that every vote cast at $\mathsf{vote}(t)$ by an honest validator with target $\preceq H$ is included in~$X$.
\end{enumerate}
Then at $\mathsf{vote}(t+1)$, every honest validator has the same $R_s$ with $R_s.h = h^\star + 1$, $(R_s, R_s.h, R_s.\slot) \in C_J$, and no honest has cast a $\just$ vote at height $> R_s.h$ at any time strictly before $\mathsf{vote}(t+1)$.
\end{lemma}

\begin{proof}
By \Cref{rev:ass:view-merge} + honest $\mathsf{syncer}(t)$ + determinism, all honest share $\Sigma^\star_t$ and $H$, with common $h_c^{(t)} = H.h = h^\star + 1$.
Hypothesis~(6) rules out any prior honest vote at $> h^\star$, so in particular none at $h^\star + 1$; (R1) fires at every honest, producing $v^{(t)}_i$ at height $h^\star + 1$ with common target $T \preceq H$ (\Cref{rev:lem:canonical-fresh-honest-vote}).
The $\geq 2n/3$ honest $\just$ votes form a $\just$-quorum for~$T$ at $h^\star + 1$.

Hypothesis~(7) + monotonicity of views places~$X$ in every honest store at $\mathsf{vote}(t+1)$; on~$X$ the quorum fires justification with $\sigma[X].J = T$ and $\sigma[X].h_j = h^\star + 1$, so $\textsf{onBlock}(X)$ inserts $(T, h^\star + 1, T.\slot) \in C_J$ into every honest store.
\Cref{rev:lem:unique-just-per-height} makes~$T$ the unique block justified at $h^\star + 1$.
Before $\mathsf{vote}(t+1)$, honest votes are at heights $\leq h^\star$ (hypothesis~(6)) or $= h^\star + 1$ (slot-$t$); under $f < n/3$ no quorum at $> h^\star + 1$ forms, so no $C_N$ tuple at $> h^\star + 1$ exists.
The argmax of $\textsf{computeRoot}$ at $\mathsf{vote}(t+1)$ therefore lands on $(T, h^\star + 1, T.\slot) \in C_J$, giving $R_s = T$ uniformly.
\end{proof}

\begin{lemma}[Healing advance across two slots]\label{rev:lem:healing-advance}
Assume:
\begin{enumerate}
  \item $f < n/3$ and honest validators follow \Cref{def:voting-rule};
  \item $\textsf{getHead}$ is deterministic;
  \item $\mathsf{start}(t) - \Delta \geq \GST$;
  \item $\mathsf{syncer}(t)$ and $\mathsf{syncer}(t+1)$ are honest; let $\Sigma^\star_t, \Sigma^\star_{t+1}$ be the common stores and $H \coloneqq \textsf{getHead}(\Sigma^\star_t)$;
  \item no honest has cast any vote at height $> H.h$ at any time strictly before $\mathsf{vote}(t)$;
  \item at $\mathsf{start}(t+1)$, $\textsf{getHead}$ applied to $\mathsf{syncer}(t+1)$'s store returns a chain $X \succeq H$ such that every vote cast at $\mathsf{vote}(t)$ by an honest validator with target $\preceq H$ is included in~$X$;
  \item by $\mathsf{start}(t+2) - \Delta$, in every honest validator's view, there exists a chain $X' \succeq X$ such that every vote cast at $\mathsf{vote}(t+1)$ by an honest validator with target $\preceq H^{(t+1)}$ is included in~$X'$, where $H^{(t+1)} \coloneqq \textsf{getHead}(\Sigma^\star_{t+1})$.
\end{enumerate}
Then at $\mathsf{vote}(t+2)$ every honest validator has the same $R_s$ with $R_s.h = H.h + 1$, $(R_s, R_s.h, R_s.\slot) \in C_J$, and no honest has cast a $\just$ vote at height $> R_s.h$ at any time strictly before $\mathsf{vote}(t+2)$.
\end{lemma}

\begin{proof}
\emph{Step 1.}
Under \Cref{rev:ass:view-merge} + honest $\mathsf{syncer}(t)$ + determinism, $h_c^{(t)} = H.h$ and honest slot-$t$ votes (either kind) all target the common $T \preceq H$; both kinds update $\mathsf{ntargets}[i]$, so on~$X$ the prefix walk at state-height $H.h$ accumulates $\geq 2n/3$ to some $B^\star \succeq T$, firing prefix notarization and inserting $(B^\star, H.h, B^\star.\slot) \in C_N$.
By \Cref{rev:ass:view-merge} + honest $\mathsf{syncer}(t+1)$ + determinism, every honest at $\mathsf{vote}(t+1)$ shares $\Sigma^\star_{t+1} = \Sigma_{\mathsf{syncer}(t+1)}$ (at $\mathsf{start}(t+1)$) and its canonical is~$X$ (hypothesis~(6)), so $\sigma[H^{(t+1)}].h = H.h + 1$; combined with \Cref{rev:lem:Rs-height-monotone} and \Cref{rev:cor:no-high-cn}, $R_s^{(t+1)}.h = H.h$.

\emph{Step 2.}
Hypothesis~(5) directly rules out any honest vote at $> H.h$ before $\mathsf{vote}(t)$; at $\mathsf{vote}(t)$ every honest casts at height~$H.h$ exactly, and no voting event occurs between $\mathsf{vote}(t)$ and $\mathsf{vote}(t+1)$, so no honest has cast at $> H.h$ before $\mathsf{vote}(t+1)$.

\emph{Step 3.}
Invoke \Cref{rev:lem:healing-one-slot} at slot~$t+1$, with $h^\star = H.h$: hypotheses (1)--(4) are inherited; (5) is Step~1; (6) is Step~2; (7) is the present hypothesis~(7).
The conclusion at $\mathsf{vote}((t+1) + 1) = \mathsf{vote}(t+2)$ is exactly what is claimed.
\end{proof}

\appendix

\section{View-merge mechanism}\label{app:view-merge}

This appendix details the view-merge mechanism abstracted by \Cref{rev:ass:view-merge} and discharges that assumption under a synchronous-slot hypothesis.
Throughout, $\Delta$ denotes the synchronous network-delivery bound.

\begin{definition}[Slot timing]\label{rev:def:slot-timing}
Each slot~$t$ is anchored at a real-valued slot time $\mathsf{start}(t)$ with $\mathsf{start}(t+1) - \mathsf{start}(t) \geq 2\Delta$ and contains three ordered events:
\begin{itemize}
  \item \emph{Freeze time} $\mathsf{start}(t) - \Delta$: every validator~$i$ records its current local store as its slot-$t$ frozen snapshot $\Sigma^{\mathrm{frozen}}_{i,t}$.
  \item \emph{Syncer time} $\mathsf{start}(t)$: $\mathsf{syncer}(t)$ broadcasts its current (unfrozen) local store to every validator; let $\Sigma^{\mathrm{sent}}_t$ be the broadcast content.
  \item \emph{Voting time} $\mathsf{start}(t) + \Delta$: every honest validator~$i$ invokes \Cref{def:voting-rule} on its slot-$t$ merged view $\Sigma^{\mathrm{merge}}_{i,t} \coloneqq \Sigma^{\mathrm{frozen}}_{i,t} \cup \Sigma^{\mathrm{sent}}_t$.
\end{itemize}
\end{definition}

\begin{definition}[Synchronous slot]\label{rev:def:synchrony-window}
Slot~$t$ is \emph{synchronous} if every message sent during $[\mathsf{start}(t) - 2\Delta, \mathsf{start}(t) + \Delta]$ is delivered to every honest validator within $\Delta$ of its sending time.
\end{definition}

\begin{lemma}[View-merge under an honest synchronous syncer]\label{rev:lem:view-merge}
If slot~$t$ is synchronous and $\mathsf{syncer}(t)$ is honest, then every honest~$i$ has $\Sigma^{\mathrm{frozen}}_{i,t} \subseteq \Sigma^{\mathrm{sent}}_t$ and $\Sigma^{\mathrm{merge}}_{i,t} = \Sigma^{\mathrm{sent}}_t$, discharging \Cref{rev:ass:view-merge} with $\Sigma^\star_t = \Sigma^{\mathrm{sent}}_t$.
\end{lemma}

\begin{proof}
Any block or vote in $\Sigma^{\mathrm{frozen}}_{i,t}$ was received by $\mathsf{start}(t) - \Delta$, hence sent by that time, hence delivered to $\mathsf{syncer}(t)$ by $\mathsf{start}(t)$ under \Cref{rev:def:synchrony-window}, so it belongs to $\Sigma^{\mathrm{sent}}_t$.
The broadcast at $\mathsf{start}(t)$ reaches every honest by $\mathsf{start}(t) + \Delta$, so $\Sigma^{\mathrm{merge}}_{i,t}$ is well-defined and equals $\Sigma^{\mathrm{sent}}_t$.
\end{proof}

\end{document}
